\documentclass[10pt]{article}
\usepackage[letterpaper,top=1in,textheight=9in,left=1.5in,textwidth=5.5in,headsep=0.25in,footskip=30pt]{geometry}
\usepackage[UTF8,scheme=plain,fontset=none,punct=quanjiao]{ctex}
\setmainfont{texgyretermes}[Extension=.otf,UprightFont=*-regular,BoldFont=*-bold,ItalicFont=*-italic,BoldItalicFont=*-bolditalic]
\setsansfont{texgyreheros}[Extension=.otf,UprightFont=*-regular,BoldFont=*-bold,ItalicFont=*-italic,BoldItalicFont=*-bolditalic]
\setmonofont{texgyrecursor}[Extension=.otf,UprightFont=*-regular,BoldFont=*-bold,ItalicFont=*-italic,BoldItalicFont=*-bolditalic]
\setCJKmainfont{Songti SC}[BoldFont={Songti SC Bold},ItalicFont={Kaiti SC}]
\setCJKsansfont{PingFang SC}
\setCJKmonofont{Songti SC}
\setCJKfamilyfont{kai}{Kaiti SC}
\newcommand{\kai}{\CJKfamily{kai}}
\linespread{1.22}
\usepackage{amsmath,amssymb}
\usepackage{graphicx,xcolor,booktabs,colortbl,multirow,array,wrapfig,caption,subcaption,enumitem}
\usepackage{fancyvrb,algorithm,algpseudocode,fancyhdr,makecell}
\usepackage[most]{tcolorbox}
\usepackage{fontawesome5}
\usepackage[round,sort]{natbib}
\setcitestyle{citesep={；},aysep={,},yysep={;}}
\usepackage{hyperref}
\hypersetup{pdftitle={RSIGame：基于递归自我改进的自主智能体游戏开发（中文译本）}}
\usepackage{url}

% ---------- 版式（仿 ICLR） ----------
\setlength{\parindent}{0pt}
\setlength{\parskip}{5.5pt plus 1pt}
\captionsetup{font=normalsize,labelsep=colon,skip=6pt}
\captionsetup[table]{position=top}
\captionsetup[sub]{font=small}
\renewcommand{\figurename}{图}
\renewcommand{\tablename}{表}
\renewcommand{\refname}{参考文献}
\renewcommand{\contentsname}{目录}
\renewcommand{\algorithmicrequire}{\textbf{输入：}}
\makeatletter
\renewcommand{\ALG@name}{算法}
\makeatother
\floatname{algorithm}{算法}
\makeatletter
\renewcommand\section{\@startsection{section}{1}{\z@}{-2.0ex plus -0.5ex minus -.2ex}{1.5ex plus 0.3ex minus0.2ex}{\large\bfseries\raggedright}}
\renewcommand\subsection{\@startsection{subsection}{2}{\z@}{-1.8ex plus -0.5ex minus -.2ex}{0.8ex plus .2ex}{\normalsize\bfseries\raggedright}}
\renewcommand\subsubsection{\@startsection{subsubsection}{3}{\z@}{-1.5ex plus -0.5ex minus -.2ex}{0.5ex plus .2ex}{\normalsize\bfseries\raggedright}}
\renewcommand\paragraph{\@startsection{paragraph}{4}{\z@}{1.2ex plus 0.5ex minus .2ex}{-1em}{\normalsize\bfseries}}
\makeatother
\setlength{\abovecaptionskip}{6pt}
\setlength{\textfloatsep}{14pt plus 2pt minus 4pt}
\renewcommand{\topfraction}{0.95}\renewcommand{\bottomfraction}{0.8}\renewcommand{\textfraction}{0.05}\renewcommand{\floatpagefraction}{0.75}
\setcounter{topnumber}{3}\setcounter{bottomnumber}{2}\setcounter{totalnumber}{5}
\def\UrlBreaks{\do\_\do\/\do\-}

\newcommand{\rsig}{RSIG\textsc{ame}}
\newcommand{\figpdf}{figs.pdf}
\newcommand{\fig}[1]{\includegraphics[page=#1]{\figpdf}}
\definecolor{rowlight}{RGB}{220,235,247}
\definecolor{rowdeep}{RGB}{143,193,227}
\definecolor{pbord}{RGB}{204,221,255}
\definecolor{pbg1}{RGB}{239,245,252}
\definecolor{pbg2}{RGB}{252,245,232}
\definecolor{pbg3}{RGB}{244,240,250}
\definecolor{pbg4}{RGB}{239,248,241}
\definecolor{pbg5}{RGB}{236,248,247}

% 提示词框
\newtcolorbox{promptbox}[2]{enhanced,breakable,colback=#1,colframe=pbord,colbacktitle=#1,coltitle=black,
  boxrule=0.5pt,arc=1.5pt,left=8pt,right=8pt,top=3pt,bottom=4pt,toptitle=3pt,bottomtitle=3pt,
  titlerule=0.4pt,titlerule style=pbord,fonttitle=\bfseries,title={#2},before skip=8pt,after skip=8pt}
\DefineVerbatimEnvironment{PV}{Verbatim}{fontsize=\small,baselinestretch=1.15}

\pagestyle{plain}
\fancypagestyle{firstpage}{\fancyhf{}\renewcommand{\headrulewidth}{0pt}
  \fancyhead[C]{\raisebox{-2pt}{\fig{1}}}\fancyfoot[C]{\thepage}}

\begin{document}
\thispagestyle{firstpage}

{\raggedright
{\LARGE\bfseries \rsig：基于递归自我改进的\\[2pt]自主智能体游戏开发\par}
\vspace{4pt}
{\large\textsc{RSIGame: Autonomous Agentic Game Development with Recursive Self-improvement}\par}
\vspace{12pt}
\hspace*{0.2in}\parbox{5in}{\linespread{1.0}\selectfont
\textbf{Wenyi Wu}$^{1*}$\quad \textbf{Minghao Fu}$^{1,2*}$\quad \textbf{Jieyu You}$^{2}$\quad \textbf{Kun Zhou}$^{1\dagger}$\quad \textbf{Siqi Liu}$^{1}$\\
\textbf{Aayush Salvi}$^{1}$\quad \textbf{Yiheng Lin}$^{2}$\quad \textbf{Ce Zhang}$^{3}$\quad \textbf{Xiaohan Lan}$^{2}$\quad \textbf{Jiahui Zhu}$^{2}$\\
\textbf{Yujie Zhong}$^{2\dagger}$\quad \textbf{Qi She}$^{2}$\quad \textbf{Biwei Huang}$^{1}$\\[2pt]
$^{1}$加州大学圣地亚哥分校\quad $^{2}$字节跳动\quad $^{3}$卡内基梅隆大学}
\par}

\vspace{14pt}
\begin{center}{\large\bfseries 摘\quad 要}\end{center}
\vspace{-4pt}
\begin{quote}
大语言模型的最新进展使自动游戏生成变得越来越可行，然而要在可玩版本的基础上可靠地进一步提升所生成游戏的质量，仍然颇具挑战。朴素的迭代式改进很容易过拟合于少量测试用例，产出脆弱的游戏：缺陷悬而未决、行为有所缺失，且难以泛化到更广泛的玩家交互。我们提出 \rsig，一个具备递归自我改进能力的自主智能体游戏开发框架。\rsig{} 将开发组织为互补的局部循环与全局循环。具体而言，局部的\emph{探索–诊断–改进}循环广泛探索可执行的游戏，对发现的问题进行诊断与优先级排序，并执行基于证据的修订，其中一份不断演化的清单持续积累新的测试与改进指引。全局循环跟踪整体质量、保留最佳检查点，并在长周期开发中检测饱和或退化。除测试时改进之外，\rsig{} 还进一步通过训练将成功的开发经验内化到生成器中。在 140 个 GameCraft-Bench 任务、两个游戏引擎和五个生成器上，\rsig{} 在相同开发预算下始终提升了游戏质量。值得注意的是，经验内化使 Qwen3.8-27B 在 Godot 上达到 61.38、在 Phaser 上达到 58.53，超过了 GPT-5.5 的单次生成得分，同时将 Qwen 的生成 token 数减少了 11 倍。
\end{quote}

\begin{center}
\href{https://github.com/WenyiWU0111/RSIGame}{\faGithub\ 代码}\qquad
\href{https://huggingface.co/RSIGame}{\raisebox{-1.5pt}{\fig{2}}\ 数据集}\qquad
\href{https://huggingface.co/spaces/RSIGame/rsigame-page}{\faGlobe\ 项目主页}
\end{center}

\begin{figure}[h]
\centering
\fig{3}
\caption{\textbf{\rsig{} 将智能体游戏开发转变为自主的递归自我改进。}它将局部的探索–诊断–改进循环与全局进度监控和控制相结合，使 Qwen3.8-27B 接近 GPT-5.5 级别的表现。}
\label{fig:teaser}
\end{figure}

\renewcommand{\thefootnote}{\fnsymbol{footnote}}
\footnotetext[1]{同等贡献。\quad $^{\dagger}$通讯作者。\quad 联系方式：\texttt{franciskunzhou@gmail.com}。}
\renewcommand{\thefootnote}{\arabic{footnote}}
\setcounter{footnote}{0}

\section{引言}
大语言模型（LLM）与视觉语言模型（VLM）的最新进展大幅扩展了自主智能体执行复杂现实任务的能力，范围涵盖从软件工程到计算机使用自动化等诸多领域\citep{yang2024,zhu2026}。自动游戏开发是另一个前景广阔的领域：构建一个完整的游戏需要协同设计游戏机制与进程，创建源代码和多模态素材，并将这些异构组件组织成一个连贯的可执行项目\citep{luo2026}。随着基础模型在长程生成、多模态理解和工具使用方面的能力迅速提升，这一设想正变得越来越现实。相应地，近期的基准测试表明，基于 VLM 的智能体已经能够把自然语言规格说明转化为具有可辨识机制和视觉内容的可玩游戏\citep{luo2026,chi2026}。

然而，游戏开发作为一个复杂的多模态工程过程，其本质决定了它依然极具挑战性。即便在成熟的人类团队中，开发过程中也常常出现意料之外的缺陷、缺失的功能以及被忽视的边界情况\citep{roque2025}。因此，一种常见的做法是采用迭代式开发流程：开发者反复测试当前构建版本、收集反馈并修订项目，直到达到期望的质量。这一过程可以自然地视为一种递归自我改进\citep{wu2024}，每一轮开发都基于新观察到的证据，在上一版本的基础上推进。受这一流程启发，我们将智能体游戏开发形式化为一个迭代循环：智能体首先生成一个初始游戏，然后反复测试并改进它，直到满足目标要求。

尽管递归自我改进天然契合智能体游戏开发，我们发现它很容易收敛到脆弱的解：大致满足了目标，却留下许多未经测试的缺陷、缺失的行为和边界情况。实际上，开发循环可能过拟合于少量测试用例，而非整体地改进游戏，导致其难以泛化到更广泛的玩家交互。为解决这一问题，我们强化了测试阶段，使其广泛探索剩余的失败与潜在改进，目标是开发出高质量的游戏，而不仅仅是可玩的游戏。具体而言，我们设计了一个\emph{探索–诊断–改进}循环：智能体首先广泛探索可执行的游戏，以发现缺陷、弱点和可增强之处；然后对发现的问题进行诊断并排列优先级，以规划下一轮改进。在整个过程中，我们维护一份\emph{演化清单}，持续积累新识别的问题以及供后续测试和修订使用的可操作指引。通过随着新证据的发现而不断扩展，这份清单持续推动开发过程超越有限的测试用例，朝着更具泛化性的整体游戏质量提升迈进。

为此，我们提出 \rsig（图~\ref{fig:overview}），一个具备递归自我改进能力的自主智能体游戏开发框架。其核心在于，\rsig{} 将开发组织为互补的局部循环与全局循环。局部循环实现了我们的探索–诊断–改进过程：反复探索可执行的游戏，诊断并排序所发现的问题，并基于积累的证据修订项目。全局循环则跨迭代监控开发过程，跟踪整体游戏质量，保留最佳检查点，并检测饱和或退化，使系统能够保持长周期的进展，而不是盲目地持续编辑。除了无需训练的自我改进之外，\rsig{} 还引入了基于训练的经验内化，将成功的开发轨迹、诊断决策和经过验证的修订转化为生成器的监督信号。这种方式使更广义的 RSI 循环能够增强底层生成器。在 140 个 GameCraft-Bench 任务、两个引擎和五个生成器上，\rsig{} 在相同开发预算下始终改进了游戏。通过迭代开发，Qwen3.8-27B 在 Godot 上达到 47.77，而 Kimi-K2.6 为 44.77\citep{moonshot2026}；经验内化将其提升至 61.38，而 Codex+GPT-5.5 的单次生成得分为 50.26\citep{openai2025,openai2026}，同时生成 token 数减少了 11 倍。在 Phaser 上，它升至 50.24，超过了 GPT-5.5 的单次生成得分 49.44。

\section{预备知识}
\paragraph{智能体游戏开发。}其目标是利用智能体实现游戏创作的自动化。给定一份自然语言的游戏规格说明，智能体将用户意图转化为一个完整且可玩的游戏。与人类游戏开发类似，智能体游戏开发必须满足两个目标：遵循用户需求，以及产出功能正确、没有缺陷的游戏。形式化地，给定游戏规格说明 $x$，我们希望学习一个生成器策略 $\pi_\theta$，用于产生游戏项目
\begin{equation}
P_0 \sim \pi_\theta(\cdot \mid x),
\end{equation}
其中 $P_0$ 既应符合预期的游戏设计，又应能正确执行。这样的项目由源代码、视觉与文本素材以及配置文件组成。因此，它的构建需要联合生成并整合代码、文本和视觉元素，使之成为一个连贯的交互系统。重要的是，游戏质量最终体现在运行时行为和玩家体验上。每个生成的游戏都包含一组可回放的演示轨迹 $\mathcal{T} = \{\tau_1, \tau_2, \ldots, \tau_N\}$，它们在评估时被回放，并依据一份隐藏的、针对具体任务的评分细则打分。

\begin{figure}[t]
\centering
\fig{4}
\caption{\rsig{} 概览。局部循环通过方向决策、主动探索、基于证据的编辑和智能体验证来自主演化游戏。全局循环监控全局质量、保留最佳检查点，并在进展饱和时引入稀疏的高层指引，从而实现渐进式的多阶段游戏演化。}
\label{fig:overview}
\end{figure}

\paragraph{我们的关注点。}由于游戏生成涉及许多紧密耦合的组件，一次性生成完全正确的项目十分困难。因此，现有工作流普遍采用基于迭代的策略：$P_0 \rightarrow P_1 \rightarrow \cdots \rightarrow P_T$。在每次迭代中，执行并测试当前项目，诊断观察到的失败，并据此修订项目。在传统游戏开发中，这一循环依赖人类开发者和试玩测试员。在本工作中，我们将同样的原则引入智能体游戏开发，但以一个完全自主的过程取代人工驱动的测试与修订：该过程递归地测试、诊断并改进其自身生成的游戏。

\section{\rsig}
我们的 \rsig{} 同时包含无需训练和基于训练的策略。无需训练的策略建立在用于迭代开发的局部循环与全局循环之上，而基于训练的策略则进一步将积累的开发经验内化到骨干游戏生成模型中。

\subsection{无需训练的递归自我改进}
我们的无需训练 RSI 策略由两部分组成：一个遵循迭代式探索–诊断–改进流程、用以识别问题并完善游戏的局部循环；以及一个监控长周期进展、保留最佳版本并检测饱和的全局循环。

\subsubsection{局部循环：探索–诊断–改进}
局部循环通过三个阶段执行细粒度的递归游戏改进：自主探索、问题诊断和迭代改进。该循环持续探索可执行的游戏，将新发现的问题和改进机会积累到一份演化清单中，并用它们指导后续修订。

\paragraph{用于自主探索的控制器与探索器。}自主探索阶段涉及两个角色互补的智能体。\emph{控制器}负责决定接下来应探索游戏的哪个方面，而\emph{探索器}与可执行的游戏交互，以收集具体的行为证据。具体来说，给定游戏规格说明 $x$、当前项目 $P_t$、开发清单 $C_t$ 以及可选的阶段级指引 $\gamma_s$，控制器提出一个开发方向 $d_t$。以此方向为条件，探索器与游戏进行有针对性的交互，并产生一条探索轨迹 $\tau_t$：
\begin{equation}
d_t \sim \pi_{\mathrm{controller}}(\cdot \mid x, P_t, C_t, \gamma_s), \qquad \tau_t \sim \pi_{\mathrm{explorer}}(\cdot \mid P_t, d_t).
\end{equation}
通过这一过程，控制器提供高层的探索意图，而探索器将其落实到观察到的游戏过程中，发现具体的缺陷、缺失的行为和潜在的改进机会，供后续诊断与修订使用。

\paragraph{用于问题诊断的编辑器与验证器。}问题诊断阶段涉及两个职责互补的智能体。\emph{编辑器}把发现的问题和交互证据转化为具体的游戏修改，而\emph{验证器}评估这些修改是否解决了预期问题且未引入退化。具体而言，以当前项目、开发方向和观察到的试玩轨迹为条件，编辑器提出一次编辑
\begin{equation}
\Delta_t \sim \pi_{\mathrm{editor}}(\cdot \mid P_t, d_t, \tau_t), \qquad P_{t+1} = \mathrm{Apply}(P_t, \Delta_t).
\end{equation}
随后，验证器与更新后的项目交互，并产生验证结果
\begin{equation}
v_t \sim \pi_{\mathrm{verifier}}(\cdot \mid P_{t+1}, d_t, \tau_t),
\end{equation}
它判定目标问题是否已被成功解决，以及该编辑是否造成了意外的副作用。验证后的结果随后被写回演化清单，为控制器规划下一轮开发提供更新的证据。

\paragraph{伴随清单更新的迭代改进。}为了在迭代之间连接探索、诊断与修订，\rsig{} 维护一份共享的开发清单 $C_t$，作为不断演化的工作状态。清单记录已发现的问题、改进机会、优先级和验证结果，使所有智能体都能基于对当前开发状态的一致视图开展工作。在每一轮中，控制器读取 $C_t$ 以确定下一个开发方向 $d_t$，探索器以新观察到的证据扩充它，编辑器处理优先级较高的条目，验证器则根据所得的游戏行为更新这些条目的状态。这些交互产生下一轮的清单
\begin{equation}
C_{t+1} = U(C_t, d_t, \tau_t, \Delta_t, v_t),
\end{equation}
它与 $P_{t+1}$ 一起传递给控制器。通过这种方式，\rsig{} 持续积累开发知识，并确保每次迭代都建立在经过验证的进展之上，而不是从头开始。

\subsubsection{全局循环：进度监控与控制}
局部循环专注于在每轮开发中改进游戏，而全局循环则跨阶段地管理整个开发过程。它由两个组件构成：跟踪累积游戏质量并保留最强检查点的\emph{游戏质量监控器}，以及决定开发应继续、终止还是在额外的高层指引下进入新阶段的\emph{进度与收敛控制}机制。

\paragraph{游戏质量监控器。}游戏质量监控器跨局部开发轮次跟踪游戏层面的进展。一个阶段从上一阶段保留的检查点开始，$P^\star_s \leftarrow P^\star_{s-1}$。在每个全局评估点，监控器将刚产生的检查点与保留的检查点进行比较，并更新
\begin{equation}
P^\star_s \leftarrow \mathrm{SelectBest}\left(P^\star_s, P_{t+1}\right).
\end{equation}
与局部验证器（评估某次具体编辑是否解决了其目标问题）不同，游戏质量监控器评估的是游戏整体的累积质量，并在整个开发过程中维护一个持久的最佳状态。重要的是，这一监控过程与基准评估器严格隔离，包括其分数、评分细则和反馈。详见附录~\ref{app:monitor}。

\paragraph{进度与收敛控制。}保留检查点 $P^\star_s$ 的演变为控制长周期开发提供了一个自然的信号。如果连续的全局评估都未能产生更好的检查点，\rsig{} 便将这种持续的无进展视为收敛。系统随后要么终止并返回 $P^\star_s$ 作为最终游戏，要么可选地引入来自人类或更强模型的新高层指引 $\gamma_{s+1}$。这类指引提供的是战略方向或创意建议，而非明确的编辑指令。于是，下一个开发阶段在 $\gamma_{s+1}$ 下从保留的最佳检查点 $P^\star_s$ 出发，使局部循环能够探索新的改进方向，同时不丢弃先前已验证的收益。重复这一过程，使 \rsig{} 能够控制长周期进展，同时保留迄今为止达到的最强解。

\subsection{基于训练的知识内化}
在无需训练的 RSI 过程中，\rsig{} 还积累了关于如何规划、实现、测试和完善游戏的可复用经验。因此，我们进一步将成功的 RSI 知识内化到骨干模型中，把自我改进从上下文优化扩展到参数优化。这一更广义的 RSI 循环改进了底层生成器，使其能够产生更强的初始生成结果，并减少高质量游戏开发所需的完善迭代次数。

\paragraph{开发经验的收集与整理。}我们收集三类监督数据：GPT-5.5\citep{openai2026}在 Godot\citep{godot2026}和 Phaser\citep{photon2026}中构建游戏的\emph{生成轨迹}；从成功生成中收集的\emph{规划轨迹}；以及使用 GLM-5.3-Flash\citep{zhipu2026}运行 \rsig{} 产生的\emph{改进轨迹}。我们只保留可执行的生成结果，以及经编辑后交互验证、且未破坏先前功能行为的改进。总计，我们获得了 2,213 条生成轨迹、2,108 条规划轨迹，以及从 4,003 个候选中得到的 2,013 个经过验证的改进轮次。

\paragraph{经验内化。}我们在整理好的规划、生成和改进轨迹上对 Qwen3.8-27B\citep{qwen2026}进行监督微调。除了从最终的游戏产物中学习之外，模型还内化了驱动成功开发的中间规划、工具使用、诊断和修订决策。这将可复用的 RSI 知识转移到模型参数中，增强了完善能力，并减少了达到高质量游戏生成所需的迭代次数。这种方式为持续 RSI 打开了大门：新获得的开发经验可以被反复内化，从而随时间推移不断改进生成器。训练细节见附录~\ref{app:train}。

\section{实验}
\subsection{实验设置}
\paragraph{基准与引擎。}我们在 GameCraft-Bench\citep{luo2026}上进行评估，它包含覆盖 15 个游戏类别的 140 个游戏开发任务。我们在 Godot\citep{godot2026}和 Phaser\citep{photon2026}中实例化相同的任务，并在两种引擎间保持任务规格、评估细则和开发协议不变。因此，Phaser 作为对 Godot 上所得结论的跨引擎稳健性评估。

\paragraph{评估协议。}遵循 GameCraft-Bench，每个生成的游戏都打包了由游戏生成器提供的可回放演示轨迹，这些轨迹被回放，并依据一份隐藏的、针对具体任务的评分细则，在机制（$M$）、深度（$D$）、视觉（$V$）和美术（$A$）四个维度上打分：
\[
Q = \textsc{Build} \times (0.15M + 0.35D + 0.15V + 0.35A),
\]
其中 $\textsc{Build} \in \{0, 1\}$ 给不可玩的游戏赋零分。我们使用 Qwen3.8-27B\citep{qwen2026}作为评审，并对三次独立的回放加打分运行取平均。隐藏的评分细则、分数和评审反馈从不暴露给开发智能体。评估稳定性和自由游玩稳健性详见附录~\ref{app:stability}和~\ref{app:freeplay}。

\paragraph{基线与开发预算。}我们在多个生成器上比较冻结的初始项目 $P_0$、Play2Code\citep{huang2026}和 \rsig。对于每个生成器，所有开发方法都从 $P_0$ 的同一份克隆出发，并使用相同的每轮工具预算。Play2Code\citep{huang2026}和 \rsig{} 每轮最多可进行 26 次工具调用，最多 30 个开发轮次。完整的模型配置、各行的具体预算、成本核算和实现细节见附录~\ref{app:config}和附录~\ref{app:cost}。

\begin{table}[tbp]
\caption{\textbf{GameCraft-Bench 上的主要结果（上：Godot；下：Phaser）。}各方法在同一骨干和预算下、从相同的基础游戏出发进行比较，结果为 140 个任务的平均；Tok. 和 Cost 为每个任务的平均计费 token 数和成本（附录~\ref{app:config}）。\textbf{粗体}标出每组内各列的最佳值；\colorbox{rowdeep}{\strut 更深}色的行内化了开发经验。}
\label{tab:main}
\centering
\small
\setlength{\tabcolsep}{4.2pt}
\begin{tabular}{@{}lccccccc@{}}
\toprule
方法 & 机制 & 深度 & 视觉 & 美术 & 总分 $\uparrow$ & Tok. $\downarrow$ & Cost $\downarrow$ \\
\midrule
\multicolumn{8}{c}{\textbf{Godot}\citep{godot2026}} \\
\midrule
\multicolumn{8}{@{}l}{\quad\kai 生成器：Codex + GPT-5.5 (high)} \\
基线（冻结的 $P_0$） & 58.9 & 51.7 & 51.4 & 44.6 & 50.26 & 0.26M & -- \\
\quad + Play2Code & 59.2 & 51.3 & 52.2 & 45.8 & 50.74 & 3.20M & \$1.06 \\
\rowcolor{rowlight}\quad + \rsig & \textbf{72.8} & \textbf{60.3} & \textbf{65.9} & \textbf{64.5} & \textbf{64.53} & \textbf{3.12M} & \textbf{\$0.88} \\
\midrule
\multicolumn{8}{@{}l}{\quad\kai 生成器：Codex + Kimi-K2.6} \\
基线（冻结的 $P_0$） & 40.2 & 30.3 & 35.4 & 21.9 & 29.63 & 3.22M & -- \\
\quad + Play2Code & 44.8 & 35.4 & 40.2 & 28.6 & 35.19 & \textbf{7.07M} & \$1.27 \\
\rowcolor{rowlight}\quad + \rsig & \textbf{51.7} & \textbf{40.1} & \textbf{47.1} & \textbf{45.4} & \textbf{44.77} & 7.39M & \textbf{\$1.18} \\
\midrule
\multicolumn{8}{@{}l}{\quad\kai 生成器：Codex + GLM-5.3-Flash} \\
基线（冻结的 $P_0$） & 35.7 & 30.2 & 32.0 & 24.9 & 29.46 & 0.44M & -- \\
\quad + Play2Code & 47.2 & 39.7 & 41.1 & 32.7 & 38.59 & \textbf{3.18M} & \textbf{\$0.99} \\
\rowcolor{rowlight}\quad + \rsig & \textbf{51.9} & \textbf{44.4} & \textbf{49.9} & \textbf{53.9} & \textbf{49.72} & 5.06M & \$1.25 \\
\midrule
\multicolumn{8}{@{}l}{\quad\kai 生成器：Codex + Qwen3.8-27B} \\
基线（冻结的 $P_0$） & 41.2 & 33.4 & 38.6 & 38.3 & 37.07 & 6.41M & -- \\
\quad + Play2Code & 44.6 & 36.1 & 42.1 & 42.5 & 40.53 & 9.92M & \$1.32 \\
\rowcolor{rowlight}\quad + \rsig & \textbf{51.2} & \textbf{40.4} & \textbf{47.5} & \textbf{53.8} & \textbf{47.77} & \textbf{9.13M} & \textbf{\$1.09} \\
\midrule
\multicolumn{8}{@{}l}{\quad\kai 生成器：Codex + Qwen3.8-27B (SFT)} \\
基线（冻结的 $P_0$） & 56.1 & 47.5 & 49.8 & 44.8 & 48.22 & 0.57M & -- \\
\quad + Play2Code & 57.4 & 48.1 & 51.6 & 47.2 & 49.71 & 3.74M & \$1.08 \\
\rowcolor{rowdeep}\quad + \rsig & \textbf{68.7} & \textbf{55.8} & \textbf{62.8} & \textbf{63.2} & \textbf{61.38} & \textbf{3.49M} & \textbf{\$0.94} \\
\midrule
\multicolumn{8}{c}{\textbf{Phaser}\citep{photon2026}} \\
\midrule
\multicolumn{8}{@{}l}{\quad\kai 生成器：OpenGame + GPT-5.5} \\
基线（冻结的 $P_0$） & 49.1 & 38.6 & 53.6 & 58.6 & 49.44 & 5.36M & -- \\
\quad + Play2Code & 57.1 & 42.0 & 60.0 & 65.2 & 55.10 & \textbf{8.34M} & \textbf{\$1.43} \\
\rowcolor{rowlight}\quad + \rsig & \textbf{59.8} & \textbf{44.2} & \textbf{62.4} & \textbf{69.7} & \textbf{58.21} & 8.38M & \$1.44 \\
\midrule
\multicolumn{8}{@{}l}{\quad\kai 生成器：OpenGame + Qwen3.8-27B} \\
基线（冻结的 $P_0$） & 40.6 & 29.5 & 43.7 & 48.7 & 40.03 & 9.03M & -- \\
\quad + Play2Code & 51.0 & 36.4 & 53.2 & 58.0 & 48.70 & 12.99M & \$1.75 \\
\rowcolor{rowlight}\quad + \rsig & \textbf{52.0} & \textbf{36.5} & \textbf{54.6} & \textbf{61.3} & \textbf{50.24} & \textbf{12.78M} & \textbf{\$1.50} \\
\midrule
\multicolumn{8}{@{}l}{\quad\kai 生成器：OpenGame + Qwen3.8-27B (SFT)} \\
基线（冻结的 $P_0$） & 47.1 & 35.4 & 48.9 & 52.4 & 45.14 & 5.42M & -- \\
\quad + Play2Code & 55.6 & 41.7 & 57.3 & 61.8 & 53.15 & 9.58M & \$1.61 \\
\rowcolor{rowdeep}\quad + \rsig & \textbf{60.4} & \textbf{46.2} & \textbf{62.3} & \textbf{68.4} & \textbf{58.53} & \textbf{9.31M} & \textbf{\$1.53} \\
\bottomrule
\end{tabular}
\end{table}

\subsection{Godot 游戏结果}
表~\ref{tab:main} 揭示了三个主要发现。第一，\emph{\rsig{} 在各生成器上都带来了大幅且一致的提升。}在全部五种设置下，它将冻结的初始游戏提高了 10.7–20.3 个总分点，提升一致地覆盖机制、深度、视觉和美术。第二，\emph{在相同计算预算下，\rsig{} 的表现显著优于其他基于递归的方法。}在相同开发预算下，\rsig{} 比 Play2Code 高出 7.2–13.8 个总分点。Play2Code 对强大的 Codex+GPT-5.5 初始化和 Qwen3.8-27B 变体几乎没有带来提升，而 \rsig{} 则大幅改进了同样的冻结游戏。这一对比表明，收益来自开发的组织方式，而不仅仅是迭代本身。最后，\emph{经验内化同时提升了质量和效率。}微调将 Qwen3.8-27B 的单次生成得分从 37.07 提高到 48.22，与 Codex+GPT-5.5 相差不到 2.0 分，同时将生成 token 数减少了 11 倍（从 6.41M 降至 0.57M）。配合 \rsig，得分进一步升至 61.38，总 token 用量降低 2.6 倍（从 9.13M 降至 3.49M）。

\begin{figure}[t]
\centering
\fig{5}
\caption{\textbf{GameCraft-Bench（40 个任务）上的开发时扩展。}在 Godot 和 Phaser 上，无论初始生成器强（GPT-5.5）或弱（Qwen3.8-27B），迭代开发都可能停滞或退化，而 \rsig{} 在各开发预算下都实现了持续改进。全局质量监控器保留迄今达到的最佳检查点，与返回最后一个检查点相比改善了质量–成本权衡。阴影区域表示 $\pm1$ 个标准误。}
\label{fig:scaling}
\vspace{6pt}
\fig{6}
\caption{\textbf{全局质量控制与经验迁移。}(\textbf{a}) 最佳检查点跟踪保护了早期收益免受后续退化的影响。(\textbf{b}) 感知饱和的停止策略以更少的开发轮次达到了相当的质量。(\textbf{c}) 内化经过验证的开发经验，在所有质量维度上改进了单次生成。}
\label{fig:global}
\end{figure}

\subsection{Phaser 游戏结果}
表~\ref{tab:main} 的 Phaser 部分表明，\rsig{} 的收益可以跨引擎迁移：它使全部三个生成器在其冻结基线之上提高了 8.8–13.4 个总分点，并在每种设置下都取得了最强的最终质量。有趣的是，Play2Code 在 Phaser 上明显强于在 Godot 上。分类细目揭示了原因：Phaser 的初始化往往机制和深度较弱、视觉和美术较强，因而总分相近，却留下了更多容易改进的功能性提升空间。Play2Code 因此可以简单地找到这些明显的缺陷并挽回可观的质量，从而缩小与 \rsig{} 的差距。尽管如此，\rsig{} 仍然始终产出最好的最终游戏，这表明即使此类改进已经占据了大部分可用的提升空间，它的优势也不止于低层次的改进。

\section{进一步分析}
\paragraph{全局质量监控能否改进开发时扩展？}图~\ref{fig:scaling} 揭示了不同开发方法随额外计算扩展时的明显差异。单纯延长预算并不能保证得到更好的游戏：Play2Code 很快陷入停滞，并且随着轮次增加经常出现退化。相比之下，\rsig{} 在 Godot 和 Phaser 上、无论初始化强弱，都能持续地把额外的开发轮次转化为更高的游戏质量。仅局部循环本身就已带来可观的收益，但其轨迹仍然波动；全局质量监控器通过保留迄今达到的最佳状态来稳定这一过程，从而在整个开发预算内实现持续改进。

\begin{figure}[t]
\centering
\fig{7}
\caption{\textbf{自适应开发跟随游戏状态并提升最终质量。}(a, b) 自适应开发根据当前瓶颈重新分配精力：构建失败后转向改进，视觉质量落后时转向美术。(c, d) 盲测成对评估一致地更偏好自适应开发，而非轮转调度。}
\label{fig:adaptive}
\end{figure}

\begin{table}[tbp]
\centering
\small
\begin{subtable}[t]{0.47\linewidth}\centering
\begin{tabular}{@{}lcc@{}}
\toprule
方法 & \makecell{有据\\精确率 $\uparrow$} & \makecell{无据目标\\/ 轮 $\downarrow$} \\
\midrule
自由评论者 & 58.6\% & 1.93 \\
智能体验证 & \textbf{72.3\%} & \textbf{0.50} \\
\bottomrule
\end{tabular}
\caption{改进前的目标溯据，40 个会话。}
\end{subtable}\hfill
\begin{subtable}[t]{0.5\linewidth}\centering
\begin{tabular}{@{}lccc@{}}
\toprule
方法 & \makecell{失败\\召回率 $\uparrow$} & \makecell{特异度\\$\uparrow$} & \makecell{平衡\\准确率 $\uparrow$} \\
\midrule
仅构建 & 0.0\% & \textbf{100.0\%} & 50.0\% \\
回放 + 验证 & \textbf{76.2\%} & 92.6\% & \textbf{84.4\%} \\
\bottomrule
\end{tabular}
\caption{改进后的失败检测，48 个经裁定的轮次。}
\end{subtable}
\captionsetup{position=bottom}
\caption{\textbf{智能体验证同时改进了目标溯据和编辑后的可靠性。}在编辑之前，它过滤掉缺乏依据的改进目标；在编辑之后，基于回放的验证能检测出仅靠编译无法发现的失败修改和退化。}
\label{tab:verify}
\end{table}

图~\ref{fig:global}(a)(b) 进一步分离出这种全局控制的益处。保留的检查点紧密跟踪各预算内的理想最佳检查点，防止后续编辑抹去早先的收益；而感知饱和的停止策略以更少的轮次达到了与长得多的固定预算运行相当的质量。综合来看，这些结果表明监控器使开发时扩展既更可靠，又更节省计算。详细的停止准则见附录~\ref{app:monitor}。

\paragraph{自适应演化重要吗？}图~\ref{fig:adaptive} 考察了局部循环自适应地选择自身开发重点（而非遵循固定的轮转调度）是否有益。自适应演化将精力投向当前的游戏状态：构建失败后，它强烈地转向改进；而一旦视觉质量成为瓶颈，它会分配明显更多的轮次给美术（图~\ref{fig:adaptive}(a)(b)）。这种自适应行为也带来了更好的最终游戏：盲测成对评估在 59\% 的比较中偏好自适应开发，而偏好轮转的只有 37\%，并且在全部四个质量维度上都具有一致的优势（图~\ref{fig:adaptive}(c)(d)）。这些结果表明，长周期开发受益于让局部循环响应游戏不断变化的需求，而不是遵循固定的日程。完整的评估细节见附录~\ref{app:adaptive}。

\paragraph{智能体验证能让局部改进更可靠吗？}表~\ref{tab:verify} 在编辑前后两个时点评估了验证的效果。基于证据的改进前验证将有据精确率从 58.6\% 提高到 72.3\%，并将每轮无依据目标数从 1.93 降至 0.50。编辑之后，基于回放的验证检测出 76.2\% 的不成功改进，平衡准确率达到 84.4\%，而仅检查构建的方法一个此类行为失败也检测不到。综合来看，这些结果表明智能体验证既提升了开发决策的质量，也提升了返回给后续轮次的反馈的可靠性。审计协议和指标定义见附录~\ref{app:audit}。

\paragraph{开发经验能否迁移到未来的生成中？}图~\ref{fig:global}(c) 表明，开发经验能够迁移回生成模型。仅在生成轨迹上训练，带来的提升在各质量维度间并不均衡；而加入规划和经过验证的改进经验后，机制、深度、视觉和美术都得到了一致的改进。提升在机制（+14.9）和深度（+14.1）上尤为显著，使总分在任何测试时开发之前就提高了 +11.1 分。这表明，规划和改进轨迹提供了可复用的知识，其价值超出了单纯模仿成功的游戏生成。

\section{相关工作}
\paragraph{游戏生成与开发基准。}近期的基准越来越多地在完整的可执行游戏而非孤立的代码片段上评估智能体。OpenGame-Bench\citep{jiang2026}通过构建有效性、视觉可用性和指令对齐来评估基于网页的游戏；GameCraft-Bench\citep{luo2026}将评估扩展到基于引擎的游戏，采用以交互为基础的机制、内容、视觉和美术度量；GameDevBench\citep{chi2026}则研究在更大代码库和素材上的多模态游戏开发。相关工作还通过基准和生成系统两个方面，进一步考察了游戏智能体中的交互以及更广泛的基于网页或代码的游戏生成\citep{lin2026,hu2024game,zhang2025,zhang2026,ma2026}。总体而言，这些工作把评估从静态的代码正确性转向可执行、可交互的游戏产物的质量。然而，它们只为单次提交的构建版本打分，至于智能体应如何在多轮开发中持续改进游戏，仍是一个悬而未决的问题——而这正是我们所研究的设定。

\paragraph{迭代式智能体游戏开发。}近期的系统将游戏生成扩展为迭代开发。Play2Code\citep{huang2026}在游玩与代码修订之间交替进行，利用每个构建版本的交互反馈来指导后续编辑；VibeGame\citep{hu2026vibe}则协调多个专门化的智能体进行持续的游戏开发。OpenGame\citep{jiang2026}进一步跨项目积累可复用的开发技能，以支持后续生成。这些工作沿袭了基于反复推理、执行与反馈的更广泛的智能体编程范式\citep{yang2024,hong2024,yao2023,shinn2023}。\rsig{} 并非引入迭代本身，而是聚焦于如何在长周期中可靠地维持迭代开发。它将自适应、基于证据的局部改进与全局的最佳检查点跟踪和饱和检测相结合。高层指引还可以在自主进展停滞后重新开启开发，从而实现多阶段的游戏演化。

\paragraph{验证与从开发经验中学习。}智能体验证已经通过多种途径得到研究，包括通用的 LLM 作为评审协议\citep{zheng2023}、智能体循环中可靠的过程控制方法\citep{wu2026}，以及通过运行时交互来验证行为的游戏专用方法\citep{jia2026,huang2026}。在 \rsig{} 中，独立验证为每一轮局部开发收尾：更新后的游戏被回放，以确认预期的改进并检测新的退化，验证结果被写回共享的开发状态。在单个任务之外，已有工作从反馈、记忆或经过验证的轨迹中学习\citep{shinn2023,madaan2023,chen2024,wang2025,pan2026,zelikman2022,zhou2026}，而游戏系统则跨生成过程复用经验\citep{jiang2026,hu2026vibe}。\rsig{} 拓展了这一思路，联合内化规划、生成和经过验证的改进经验，以增强未来的游戏生成与开发。

\section{结论}
我们提出了 \rsig，一个将游戏创作形式化为递归自我改进过程的自主智能体游戏开发框架。\rsig{} 通过广泛的探索、结构化的诊断、按优先级的改进以及一份不断扩展已识别问题和改进机会集合的演化清单，强化了开发循环。其局部的探索–诊断–改进循环支持对可执行游戏进行基于证据的修订，而全局循环跟踪长周期进展、保留最佳检查点，并检测饱和或退化。我们还进一步引入了基于训练的经验内化，将成功的开发轨迹迁移回骨干游戏生成模型，使未来的游戏创作获得超越单纯测试时完善的改进。在 140 个 GameCraft-Bench 任务、两个游戏引擎和五个生成器上的实验表明，\rsig{} 在相同开发预算下始终提升游戏质量，使较小的开放模型在迭代开发之后能够达到甚至超越明显更强的生成器。

\section*{AI 使用声明}
大语言模型是本工作研究方法中不可或缺的一部分。GPT-5.5 用于游戏生成和规划数据的收集；GLM-5.3-Flash 用于游戏改进、验证以及改进数据的收集；Qwen3.8-27B 用作游戏测试智能体、基准评估评审，以及经验内化的基础模型。此外，Claude Opus 5 被用于独立的验证审计。

在研究和稿件准备过程中，我们还使用了基于 LLM 的助手进行语言编辑、呈现润色和编程辅助。实验由作者设计，作者审阅了所得的产物和测量结果，并核实了所有报告的结果与科学论断。作者对本文内容承担全部责任。

\section*{伦理声明}
我们的研究关注通用的游戏生成与游戏开发任务。我们在生成过程中施加了安全约束，并从评估和数据整理中排除了包含不安全或敏感内容的产物。实验不使用私人用户数据、个人信息或用户生成内容。

人工参与仅限于游戏评估和高层开发指引。我们研究的是由此得到的游戏评估和指引，而非参与者的属性，且不收集任何敏感个人信息。发布的数据集和产物经过筛查，移除了机器特定路径、凭据和个人信息。

\section*{可复现性声明}
所有报告的结果都可以追溯到保留的实验产物和有文档记录的协议。表~\ref{tab:repro} 为发布的产物以及说明实验流程各组成部分的附录章节提供了统一索引。

\textbf{受控比较。}对于每项比较，初始项目 $P_0$ 都是冻结的，每种开发方法都获得同一份克隆、相同的改进骨干和相同的每轮工具预算。各行具体的开发预算在附录~\ref{app:config} 中明确报告。每个产出的游戏都被独立回放并打分三次，而基准评分细则、分数和评审输出对所有开发智能体保持隐藏。回放和评审的变异性在附录~\ref{app:stability} 中量化。

\textbf{产物可得性。}完整的实现、提示词和评估流程可在我们的代码仓库\footnote{\url{https://github.com/WenyiWU0111/RSIGame}}中获取，具有代表性的演化游戏可在项目主页\footnote{\url{https://huggingface.co/spaces/RSIGame/rsigame-page}}上游玩。冻结的基础项目、训练轨迹和逐轮运行树正在内部审核中，获批后将发布。我们还在 Hugging Face\footnote{\url{https://huggingface.co/datasets/RSIGame/RSIGame-TableArtifacts}}上发布了支撑所报告结果的完整评估产物，包括 51,644 个评分文件，含逐任务的分数、评审输出和回放报告。它们的发布状态及相应的协议文档汇总于表~\ref{tab:repro}。

\clearpage
\begin{thebibliography}{40}
\small
\setlength{\itemsep}{2pt}

\bibitem[Anthropic(2026)]{anthropic2026}
Anthropic.
\newblock Claude Opus 5 system card.
\newblock \url{https://www.anthropic.com/claude-opus-5-system-card}, July 2026.

\bibitem[Chen et~al.(2024)]{chen2024}
Xinyun Chen, Maxwell Lin, Nathanael Schärli, and Denny Zhou.
\newblock Teaching large language models to self-debug.
\newblock In \emph{International Conference on Learning Representations (ICLR)}, 2024. arXiv:2304.05128.

\bibitem[Chi et~al.(2026)]{chi2026}
Wayne Chi, Yixiong Fang, Arnav Yayavaram, Siddharth Yayavaram, Seth Karten, Qiuhong Anna Wei, Runkun Chen, Alexander Wang, Valerie Chen, Ameet Talwalkar, and Chris Donahue.
\newblock Gamedevbench: Evaluating agentic capabilities through game development.
\newblock In \emph{Proceedings of the 43rd International Conference on Machine Learning (ICML)}, 2026. arXiv:2602.11103.

\bibitem[Dettmers et~al.(2023)]{dettmers2023}
Tim Dettmers, Artidoro Pagnoni, Ari Holtzman, and Luke Zettlemoyer.
\newblock Qlora: Efficient finetuning of quantized llms.
\newblock In \emph{Advances in Neural Information Processing Systems (NeurIPS)}, 2023. arXiv:2305.14314.

\bibitem[{Godot Engine contributors}(2026)]{godot2026}
Godot Engine contributors.
\newblock Godot engine.
\newblock \url{https://godotengine.org}, 2026.

\bibitem[Hong et~al.(2024)]{hong2024}
Sirui Hong, Mingchen Zhuge, Jiaqi Chen, Xiawu Zheng, Yuheng Cheng, Ceyao Zhang, Jinlin Wang, Zili Wang, Steven Ka~Shing Yau, Zijuan Lin, Liyang Zhou, Chenyu Ran, Lingfeng Xiao, Chenglin Wu, and Jürgen Schmidhuber.
\newblock Metagpt: Meta programming for a multi-agent collaborative framework.
\newblock In \emph{International Conference on Learning Representations (ICLR)}, 2024. arXiv:2308.00352.

\bibitem[Hu et~al.(2024)]{hu2024game}
Chengpeng Hu, Yunlong Zhao, and Jialin Liu.
\newblock Game generation via large language models.
\newblock In \emph{2024 IEEE Conference on Games (CoG)}, 2024.

\bibitem[Hu et~al.(2022)]{hu2022lora}
Edward~J. Hu, Yelong Shen, Phillip Wallis, Zeyuan Allen-Zhu, Yuanzhi Li, Shean Wang, Lu~Wang, and Weizhu Chen.
\newblock Lora: Low-rank adaptation of large language models.
\newblock In \emph{International Conference on Learning Representations (ICLR)}, 2022. arXiv:2106.09685.

\bibitem[Hu et~al.(2026)]{hu2026vibe}
Wenbo Hu, Ken Li, Jiazhe Wei, Yukang Cao, Weiyi Hong, Jiayi Dai, Chenjun Bai, Jiajun Liang, Yucheng Liao, Ruichuan An, Zeyu Lou, Haofan Wang, Yueming Lyu, Ziwei Liu, and Chenyang Si.
\newblock Vibegame: Prompt-to-game development with ai-native engine and self-evolving adversarial agent team.
\newblock Technical report, 2026. \url{https://vibegame.tettet.org/technical_report.pdf}.

\bibitem[Huang et~al.(2026)]{huang2026}
Yixu Huang, Bo~Li, Na~Li, Zhe Wang, Kaijie Chen, Haonan Ge, Qingyi Si, Yuanzhe Shen, Ruihan Yang, Guangjing Wang, and Hongcheng Guo.
\newblock Gui agents for continual game generation, 2026.

\bibitem[Jacobs et~al.(2023)]{jacobs2023}
Sam~Ade Jacobs, Masahiro Tanaka, Chengming Zhang, Minjia Zhang, Shuaiwen~Leon Song, Samyam Rajbhandari, and Yuxiong He.
\newblock Deepspeed ulysses: System optimizations for enabling training of extreme long sequence transformer models, 2023.

\bibitem[Jia et~al.(2026)]{jia2026}
Chaobo Jia, Ruipeng Wan, Ting Sun, Weihao Tan, Borui Wan, Yuxuan Tong, Guangming Sheng, and Hong Xu.
\newblock Gamegen-verifier: Parallel keypoint-based verification for llm-generated games via runtime state injection, 2026.

\bibitem[Jiang et~al.(2026)]{jiang2026}
Yilei Jiang, Jinyuan Hu, Qianyin Xiao, Yaozhi Zheng, Ruize Ma, Kaituo Feng, Jiaming Han, Tianshuo Peng, Kaixuan Fan, Manyuan Zhang, and Xiangyu Yue.
\newblock Opengame: Open agentic coding for games, 2026.

\bibitem[Lin et~al.(2026)]{lin2026}
Mingxian Lin, Shengju Qian, Yuqi Liu, Yi-Hua Huang, Yiyu Wang, Wei Huang, Yitang Li, Fan Zhang, Zeyu Hu, Lingting Zhu, Xin Wang, and Xiaojuan Qi.
\newblock Omnigamearena: A unified ue5 benchmark for vlm game agents with improvement dynamics, 2026.

\bibitem[Loshchilov \& Hutter(2019)]{loshchilov2019}
Ilya Loshchilov and Frank Hutter.
\newblock Decoupled weight decay regularization.
\newblock In \emph{International Conference on Learning Representations (ICLR)}, 2019. arXiv:1711.05101.

\bibitem[Luo et~al.(2026)]{luo2026}
Tongxu Luo, Rongsheng Wang, Jiaxi Bi, Chenming Xu, Zhengyang Tang, Jianlong Chen, Juhao Liang, Ke~Ji, Shuqi Guo, Yuhao Du, Fan Bu, Wenyu Du, Xiaotong Zhang, Kyle Li, Shaobo Wang, Linfeng Zhang, Yuxuan Liu, Xin Lai, Chenxin Li, Yiduo Guo, Zhexin Zhang, Xinyuan Wang, Tianyi Bai, Ziniu Li, and Benyou Wang.
\newblock Gamecraft-bench: Can agents build playable games end-to-end in a real game engine?, 2026.

\bibitem[Ma et~al.(2026)]{ma2026}
Hongnan Ma, Han Wang, Shenglin Wang, Tieyue Yin, Yiwei Shi, Yucong Huang, Yingtian Zou, Muning Wen, and Mengyue Yang.
\newblock Creativegame: Toward mechanic-aware creative game generation, 2026.

\bibitem[Madaan et~al.(2023)]{madaan2023}
Aman Madaan, Niket Tandon, Prakhar Gupta, Skyler Hallinan, Luyu Gao, Sarah Wiegreffe, Uri Alon, Nouha Dziri, Shrimai Prabhumoye, Yiming Yang, Shashank Gupta, Bodhisattwa~Prasad Majumder, Katherine Hermann, Sean Welleck, Amir Yazdanbakhsh, and Peter Clark.
\newblock Self-refine: Iterative refinement with self-feedback.
\newblock In \emph{Advances in Neural Information Processing Systems (NeurIPS)}, 2023. arXiv:2303.17651.

\bibitem[{Moonshot AI}(2026)]{moonshot2026}
Moonshot AI.
\newblock Kimi-K2.6.
\newblock Model card, Hugging Face, April 2026. \url{https://huggingface.co/moonshotai/Kimi-K2.6}.

\bibitem[OpenAI(2025)]{openai2025}
OpenAI.
\newblock Introducing Codex.
\newblock \url{https://openai.com/index/introducing-codex/}, May 2025.

\bibitem[OpenAI(2026)]{openai2026}
OpenAI.
\newblock GPT-5.5 system card.
\newblock \url{https://openai.com/index/gpt-5-5-system-card/}, April 2026.

\bibitem[Pan et~al.(2026)]{pan2026}
Wenbo Pan, Shujie Liu, Xiangyang Zhou, Shiwei Zhang, Wanlu Shi, Mirror Xu, and Xiaohua Jia.
\newblock M$^\star$: Every task deserves its own memory harness, 2026.

\bibitem[{Photon Storm}(2026)]{photon2026}
Photon Storm.
\newblock Phaser: A fast, free and fun open source html5 game framework.
\newblock \url{https://phaser.io}, 2026.

\bibitem[{Qwen Team}(2026)]{qwen2026}
Qwen Team.
\newblock Qwen3.8-27B.
\newblock Model card, Hugging Face, August 2026. \url{https://huggingface.co/Qwen/Qwen3.8-27B}.

\bibitem[Roque et~al.(2025)]{roque2025}
Alejandro Roque, Juan~P. Sotomayor, Dionny Santiago, and Peter~J. Clarke.
\newblock A literature review of software testing practices and frameworks in the video gaming industry.
\newblock \emph{Software Testing, Verification and Reliability}, 35(2):e70001, 2025. doi: 10.1002/stvr.70001.

\bibitem[Shinn et~al.(2023)]{shinn2023}
Noah Shinn, Federico Cassano, Edward Berman, Ashwin Gopinath, Karthik Narasimhan, and Shunyu Yao.
\newblock Reflexion: Language agents with verbal reinforcement learning.
\newblock In \emph{Advances in Neural Information Processing Systems (NeurIPS)}, 2023. arXiv:2303.11366.

\bibitem[Wang et~al.(2025)]{wang2025}
Yinjie Wang, Ling Yang, Ye~Tian, Ke~Shen, and Mengdi Wang.
\newblock Co-evolving llm coder and unit tester via reinforcement learning.
\newblock In \emph{Advances in Neural Information Processing Systems (NeurIPS)}, 2025. Spotlight. arXiv:2506.03136.

\bibitem[Wu et~al.(2024)]{wu2024}
Tianhao Wu, Weizhe Yuan, Olga Golovneva, Jing Xu, Yuandong Tian, Jiantao Jiao, Jason Weston, and Sainbayar Sukhbaatar.
\newblock Meta-rewarding language models: Self-improving alignment with LLM-as-a-meta-judge, 2024.
\newblock \url{https://arxiv.org/abs/2407.19594}.

\bibitem[Wu et~al.(2026)]{wu2026}
Wenyi Wu, Sibo Zhu, Kun Zhou, Aayush Salvi, Zixuan Song, and Biwei Huang.
\newblock Structagent: Harness long-horizon digital agents with unified causal structure, 2026.
\newblock \url{https://arxiv.org/abs/2607.11388}.

\bibitem[Yang et~al.(2024)]{yang2024}
John Yang, Carlos~E. Jimenez, Alexander Wettig, Kilian Lieret, Shunyu Yao, Karthik Narasimhan, and Ofir Press.
\newblock Swe-agent: Agent-computer interfaces enable automated software engineering.
\newblock In \emph{Advances in Neural Information Processing Systems (NeurIPS)}, 2024. arXiv:2405.15793.

\bibitem[Yao et~al.(2023)]{yao2023}
Shunyu Yao, Jeffrey Zhao, Dian Yu, Nan Du, Izhak Shafran, Karthik Narasimhan, and Yuan Cao.
\newblock React: Synergizing reasoning and acting in language models.
\newblock In \emph{International Conference on Learning Representations (ICLR)}, 2023. arXiv:2210.03629.

\bibitem[Zelikman et~al.(2022)]{zelikman2022}
Eric Zelikman, Yuhuai Wu, Jesse Mu, and Noah~D. Goodman.
\newblock Star: Bootstrapping reasoning with reasoning.
\newblock In \emph{Advances in Neural Information Processing Systems (NeurIPS)}, 2022. arXiv:2203.14465.

\bibitem[Zhang et~al.(2025)]{zhang2025}
Wei Zhang, Jack Yang, Renshuai Tao, Lingzheng Chai, Shawn Guo, Jiajun Wu, Xiaoming Chen, Ganqu Cui, Ning Ding, Xander Xu, Hu~Wei, and Bowen Zhou.
\newblock V-gamegym: Visual game generation for code large language models, 2025.

\bibitem[Zhang et~al.(2026)]{zhang2026}
Wenyu Zhang, Guoliang You, Tianlun, Haotian Zhao, Tianshu Zhu, Haoran Wang, Xiaoxuan Tang, Mingyang Dai, Jingnan Gu, Daxiang Dong, and Jianmin Wu.
\newblock Webgamebench: Requirement-to-application evaluation for coding agents via browser-native games, 2026.

\bibitem[Zheng et~al.(2023)]{zheng2023}
Lianmin Zheng, Wei-Lin Chiang, Ying Sheng, Siyuan Zhuang, Zhanghao Wu, Yonghao Zhuang, Zi~Lin, Zhuohan Li, Dacheng Li, Eric~P. Xing, Hao Zhang, Joseph~E. Gonzalez, and Ion Stoica.
\newblock Judging llm-as-a-judge with mt-bench and chatbot arena.
\newblock In \emph{Advances in Neural Information Processing Systems (NeurIPS), Datasets and Benchmarks Track}, 2023. arXiv:2306.05685.

\bibitem[{Zhipu AI}(2026)]{zhipu2026}
Zhipu AI.
\newblock GLM-5.3-Flash.
\newblock Model card, Hugging Face, August 2026. \url{https://huggingface.co/zai-org/GLM-5.3-Flash}.

\bibitem[Zhou et~al.(2026)]{zhou2026}
Chenyu Zhou, Qiliang Jiang, Shuning Wu, and Xu~Zhou.
\newblock The verifier is the curriculum: Precision sets the return on search in code self-distillation, 2026.

\bibitem[Zhu et~al.(2026)]{zhu2026}
Sibo Zhu, Shicheng Fan, Xinyue Wang, Wenyi Wu, Kun Zhou, and Biwei Huang.
\newblock Rsiagent: Autonomous exploration for recursive self-improvement in new environments, 2026.
\newblock \url{https://arxiv.org/abs/2609.15364}.

\end{thebibliography}

\clearpage
\appendix
\renewcommand{\thesection}{\Alph{section}}
\begin{center}{\Large\bfseries 附录概览}\end{center}
本附录提供支撑正文的实现细节、实验协议、训练规格和扩展分析。

\vspace{4pt}
{\small
\newcommand{\tocA}[3]{\textbf{#1}\quad\textbf{#2}\dotfill\pageref{#3}\par\vspace{-1pt}}
\newcommand{\tocD}[1]{\hspace*{1.6em}{\color{black!65}#1}\par\vspace{2pt}}
\newcommand{\tocB}[3]{\hspace*{1.6em}#1\quad #2\dotfill\pageref{#3}\par\vspace{-1pt}}
\tocA{A}{可复现性索引}{app:repro}
\tocD{产物、发布状态，以及各项协议细节的说明位置。}
\tocA{B}{局限性}{app:limit}
\tocD{自主流程的适用范围、所评估的维度以及项目规模。}
\tocA{C}{\rsig{} 的系统设计与实现}{app:system}
\tocD{递归开发系统及其智能体接口的完整规格。}
\tocB{C.1}{递归开发流程}{app:proc}
\tocB{C.2}{全局质量监控}{app:monitor}
\tocB{C.3}{高层导演界面}{app:director}
\tocB{C.4}{智能体提示词与输出约定}{app:prompts}
\vspace{3pt}
\tocA{D}{实验协议与评估可靠性}{app:exp}
\tocD{所报告实验背后的配置、核算与可靠性协议。}
\tocB{D.1}{实验配置}{app:config}
\tocB{D.2}{Token 与成本核算}{app:cost}
\tocB{D.3}{回放与评分稳定性}{app:stability}
\tocB{D.4}{自由游玩评估}{app:freeplay}
\tocB{D.5}{主要结果的统计可靠性}{app:stat}
\tocB{D.6}{自适应方向策略评估}{app:adaptive}
\tocB{D.7}{智能体验证审计}{app:audit}
\tocB{D.8}{多阶段指引评估}{app:multistage}
\vspace{3pt}
\tocA{E}{训练数据与经验内化}{app:train}
\tocD{游戏开发经验的构建、整理与内化。}
\tocB{E.1}{训练语料的构建与整理}{app:corpus}
\tocB{E.2}{微调配置}{app:ft}
\tocB{E.3}{训练数据消融}{app:ablation}
\vspace{3pt}
\tocA{F}{扩展评估与案例研究}{app:extended}
\tocD{扩展比较、定性轨迹以及按任务类别的细分结果。}
\tocB{F.1}{与多智能体“生成–验证”系统的比较}{app:vibegame}
\tocB{F.2}{定性演化案例研究}{app:cases}
\tocB{F.3}{按类别的结果}{app:family}
}
\clearpage

\section{可复现性索引}\label{app:repro}
表~\ref{tab:repro} 即可复现性声明中提到的统一索引：对于每项产物，它给出该产物包含的内容、是否已发布，以及相应协议或实现细节的说明位置。

\begin{table}[htbp]
\caption{\textbf{可复现性索引。}产物、发布状态，以及相应实现与协议细节的位置。}
\label{tab:repro}
\centering\small
\begin{tabular}{@{}p{0.27\linewidth}p{0.47\linewidth}p{0.2\linewidth}@{}}
\toprule
条目 & 内容 / 规格 & 位置 \\
\midrule
\multicolumn{3}{@{}l}{\textit{\kai 已发布的产物}} \\
实现 & 局部/全局循环、全局质量监控器、导演评审工具和评分框架 & 代码，已上线 \\
提示词与输出约定 & 所有智能体使用的完整提示词和结构化输出 & 代码，已上线；附录~\ref{app:prompts} \\
可游玩的演化过程 & 六个具代表性的游戏及其保留的检查点 & 项目主页，已上线 \\
冻结的基础项目 $P_0$ & 每个任务、每个生成器的起始项目 & 待发布 \\
评估产物 & 逐任务的细则分数、评审输出和回放报告（51,644 个文件） & HF，已上线$^\dagger$ \\
训练语料 & 2,213 条生成轨迹、2,108 份规划和 2,013 个经验证的改进轮次 & 待发布 \\
运行树 & 长周期开发运行的逐轮项目快照 & 待发布 \\
\midrule
\multicolumn{3}{@{}l}{\textit{\kai 协议与实现细节}} \\
递归开发循环 & 一个局部轮次、全局检查点更新和停止逻辑 & 附录~\ref{app:proc} \\
全局质量监控器 & 最佳检查点选择、饱和准则和耐心值 $K$ & 附录~\ref{app:monitor} \\
模型分配 & 用于生成、探索、编辑、验证和评审的模型 & 表~\ref{tab:models} \\
超参数 & 循环、监控器和评分参数 & 表~\ref{tab:params} \\
开发预算 & 每个报告行的轮次与工具调用预算 & 附录~\ref{app:config} \\
Token 与成本核算 & 计费 token、排除的成本和定价 & 附录~\ref{app:cost} \\
评分稳定性 & 回放与评审的变异性 & 附录~\ref{app:stability} \\
语料构建 & 抽取、验证过滤和污染检查 & 附录~\ref{app:corpus} \\
微调 & 训练数据混合与超参数 & 附录~\ref{app:ft} \\
导演评审 & 导演的输入、交互预算和指引格式 & 附录~\ref{app:director} \\
按类别的结果 & 按任务类别细分的主表结果 & 附录~\ref{app:family} \\
\bottomrule
\end{tabular}

\vspace{3pt}
\parbox{0.92\linewidth}{\footnotesize $^\dagger$ 发布于 Hugging Face：\texttt{RSIGame/RSIGame-TableArtifacts}。发布前已移除机器特定路径和可识别身份的术语，导出的文件经过重新解析和重新扫描。}
\end{table}

\section{局限性}\label{app:limit}
\rsig{} 的核心流程是完全自主的：自适应的局部开发、全局质量监控、最佳检查点跟踪和感知饱和的停止都不需要人工干预。我们另外研究了一种可选的有引导扩展：只有在自主改进饱和之后，才引入来自人类或更强模型的稀疏高层指引。目前的研究只在有限的多阶段开发案例上评估了这一机制；关于何时应调用指引、多少指引是有益的、以及人类指引与模型指引有何差异，更广泛的研究留待未来工作。

我们的评估也集中于机制、深度、视觉和美术四个维度。这些维度并不能完全刻画更高层次的属性，例如原创性、叙事质量、长期玩家参与度或主观乐趣。

最后，我们的实验关注的是相对紧凑、能够在实际的智能体预算内迭代开发的游戏。将自主递归开发扩展到规模大得多、周期更长、素材更丰富且跨系统依赖更复杂的项目，仍然是一个重要的方向。

\section{\rsig{} 的系统设计与实现}\label{app:system}
\subsection{递归开发流程}\label{app:proc}
算法~\ref{alg:rsig} 是第~3 节循环的完整写法：内循环的一个轮次、全局质量监控器运行的检查点，以及饱和可能导致的两种结果——返回冠军版本，或接受一条指引并开启下一阶段。

多阶段结构由计数器 $c$ 和阶段索引 $s$ 体现。一个检查点要么晋升出新的冠军（将 $c$ 归零），要么没有找到更好的版本（$c$ 加一）；连续 $K$ 个检查点都没有晋升，就是本文所说的饱和。饱和本身并不结束开发：它结束的是当前阶段。如果有可用的指引，就针对冠军写一份简报，下一阶段从该冠军（而不是上一轮的构建）重新开始，$c$ 被重置，循环在 $\gamma_{s+1}$ 引导后续每一轮的情况下继续。只有当某个阶段饱和且不再给出任何指引时，开发才会结束。本文的运行使用 $E = 3$ 和 $K = 3$。

\begin{algorithm}[h]
\caption{\rsig{} 递归游戏开发}
\label{alg:rsig}
\small
\begin{algorithmic}[1]
\Require 游戏规格说明 $x$，初始项目 $P_0$，初始清单 $C_0$，检查点间隔 $E$，耐心值 $K$
\State $P^\star_0 \gets P_0,\ \gamma_0 \gets \varnothing,\ t \gets 0,\ s \gets 0,\ c \gets 0$ \Comment{冠军、指引、轮次、阶段、空闲检查点数}
\While{仍有开发预算}
  \State $d_t \sim \pi_{\mathrm{controller}}(\cdot \mid x, P_t, C_t, \gamma_s)$ \Comment{清单和指引引导本轮}
  \State $\tau_t \sim \pi_{\mathrm{explorer}}(\cdot \mid P_t, d_t)$
  \State $\Delta_t \sim \pi_{\mathrm{editor}}(\cdot \mid P_t, d_t, \tau_t),\ P_{t+1} \gets \mathrm{Apply}(P_t, \Delta_t)$
  \State $v_t \sim \pi_{\mathrm{verifier}}(\cdot \mid P_{t+1}, d_t, \tau_t)$ \Comment{改进是否可观察，什么发生了退化}
  \State $C_{t+1} \gets U(C_t, d_t, \tau_t, \Delta_t, v_t)$ \Comment{下一轮赖以构建的工作记忆}
  \If{$t \bmod E = 0$} \Comment{检查点：由监控器比较，而非循环本身}
    \State $P' \gets \mathrm{SelectBest}(P^\star_s, P_{t+1})$
    \If{$P' \neq P^\star_s$}
      \State $P^\star_s \gets P';\ c \gets 0$ \Comment{新冠军；本阶段仍有回报}
    \Else
      \State $c \gets c + 1$ \Comment{又一个没有更好版本的检查点}
    \EndIf
    \If{$c \geq K$} \Comment{饱和：$K$ 个检查点无新冠军}
      \If{没有可用的指引} \Return $P^\star_s$ \EndIf
      \State $\gamma_{s+1} \gets \textsc{Guidance}(x, P^\star_s)$ \Comment{针对冠军的简报开启下一阶段}
      \State $P^\star_{s+1} \gets P^\star_s,\ P_{t+1} \gets P^\star_s,\ s \gets s+1,\ c \gets 0$ \Comment{下一阶段从这里重新开始}
    \EndIf
  \EndIf
  \State $t \gets t + 1$
\EndWhile
\State \Return $P^\star_s$
\end{algorithmic}
\end{algorithm}

\subsection{全局质量监控}\label{app:monitor}
\paragraph{比较什么。}监控器从不单独给某个版本打分。在每个检查点，它在同一组演示上回放保留的冠军 $P^\star_s$ 和候选构建，逐个场景地将两段录像配对，并且只问两者之间发生了什么变化。比较器不会被告知候选来自哪一轮、该改进试图实现什么，也不会被告知已花费了多少预算：方向和预算在上一层决定，如果在这里提供它们，开发意图就可能影响模型所报告看到的内容。

\paragraph{量表。}代理评分细则的每条准则都在三个锚定等级上读取——1.0 表示按描述工作，0.5 表示部分工作、间歇工作或仅在某些情况下工作，0.0 表示不工作——冠军和候选分别读取，并且只针对两段录像都展示到的准则，从而使差异反映的是游戏本身的差异，而不是两次回放恰好到达的内容不同。各准则读数被合并为 $[0,1]$ 区间内的加权总值，并分为两半报告：来自当前阶段简报的准则，以及来自任务文档的准则。

\paragraph{采纳。}当简报那一半提高至少 0.05、且任务那一半下降不超过 0.05 时，候选取代冠军；否则保留冠军。任务那一半起到的是底线而非目标的作用，因此只要构建没有让任务已要求的内容退化，就可以因为推进了本阶段自身的目标而被采纳。如果没有阶段准则，该规则就退化为总值提高 0.05。

\paragraph{检查点与停止。}检查点每三轮进行一次，这是能够包含一次改进、其验证和一次回放的最短间隔，因此也是能观察到差异的最短窗口。一旦冠军连续 $K$ 个检查点保持不变，开发即停止；我们使用 $K = 3$。图~\ref{fig:global}(b) 在相同的运行上将 $K = 1 \ldots 4$ 与固定预算进行比较：较小的 $K$ 停得更早、交付得更少，较大的 $K$ 以固定预算轮次的一小部分接近其得分，而 $K = 3$ 恰好位于拐点。

\paragraph{与评估器隔离。}监控器与基准评审只共享回放机制。监控器读取的是一份单独的代理评分细则，从不读取基准的保留细则，也从不读取其分数或反馈；该模块完全不导入评审。这种隔离在内部同样成立：比较器是唯一看到像素的组件，而决定方向和停止的阶段读取的是它的结构化计数而非画面帧，因此没有任何一层既能解读图像，又能依据自己的解读采取行动。

\subsection{高层导演界面}\label{app:director}
第~3 节描述了阶段简报对循环的作用。本附录描述简报从何而来。本文报告的每一份简报——无论来自人类还是模型——都是在下述界面中、按同一协议写成的，因此“由人指导这次运行”和“由模型指导这次运行”之间，除了导演本身之外别无差异。

\subsubsection{导演获得什么}
导演评审的是某个游戏的某一个构建版本。图~\ref{fig:director} 就是该页面：游戏在左侧面板中实时运行，框架获得焦点时接管键盘；右栏显示预算、生成该游戏所依据的任务规格说明，以及已发生开发过程的摘要；下方的表单就是简报。

这份摘要刻意保持简略。它报告已运行了多少个自主轮次、其中哪些在全局质量监控器的每三轮检查点上成为冠军、正在评审的是哪个构建，以及这些轮次如何在功能性工作和呈现性工作之间分配。它不报告任何构建在任何一轮的任何分数。

这种省略正是关键所在。导演被问到的是游戏接下来需要什么，而分数回答的是另一个问题——上一轮是否奏效——而且用的是循环本身已经在优化的那种“货币”。关于运行的其他任何信息也都不会出现在页面上：没有源代码、没有日志、没有改进过程记录、没有其他导演的评审。这个工具展示的只是一个可玩的游戏，以及它本应成为的那份文档。

\begin{figure}[t]
\centering
\fig{8}
\caption{导演评审页面（会话进行中），展示的是已经历六个监控器检查点的 \textit{Alley Brawlers} 构建：实时游戏、预算、任务规格说明和开发摘要——冠军出现在 R03 和 R09，此后再无更新，以及十八轮在功能性与呈现性工作之间的分配。画面是一个真实的游戏瞬间，展示了两位导演随后用不同措辞描述的缺陷：玩家的格斗者已经散成一条杂乱的精灵图块带。其下的简报表单见图~\ref{fig:brief}。}
\label{fig:director}
\end{figure}

\subsubsection{同一工具，同一预算}
导演有 60 次操作和 2 次重置。按一次键无论按住多久都算一次操作，点击一次也算一次，而观看和等待是免费的，因此预算买的是交互而不是耐心。图~\ref{fig:director} 中的计数器就是会话实际执行的计数器，每位导演实际花费的次数随简报一起保存。

人类通过浏览器游玩。模型则通过一个命令行客户端游玩——\texttt{open}、\texttt{press}、\texttt{hold}、\texttt{click}、\texttt{wait}、\texttt{look}、\texttt{reset}、\texttt{state}、\texttt{submit}——它驱动同一服务器上的同一个会话对象，通过同样的计数器花费同样的预算，并在每次操作后返回捕获的画面帧供模型阅读。两位导演都看不到对方的传输方式；二者面对的是同一个正在运行的游戏。

两位导演评审的是完全相同的字节。每份简报都记录了其所针对的项目树的哈希值，对于 \textit{Alley Brawlers}，两者读到的都是 \texttt{c12fe90b30b6c176}。人类导演用满了 60 次操作，没有重置；模型用了 25 次，没有重置。

\subsubsection{模型如何担任导演}
模型导演获得一份冻结的、带版本号的提示词（\texttt{director-prompt/2}），在一次运行中保持不变；修改它就构成一次新的运行。该提示词让模型扮演一位资深游戏导演，评审一个自主开发已陷入停滞的构建；要求它决定下一阶段应该是什么，而不是罗列缺陷；并禁止给出实现层面的指令。

有两条规则承担了方法论上的作用。模型对游戏的一切了解都必须来自会话本身：读取、列出或搜索任何其他文件——任何源代码、日志、分数或其他评审——都会使该评审作废。并且简报必须通过表单提交所用的同一个端点、以同一种格式提交，且只能提交一次；提交即结束会话。

\subsubsection{简报是什么}
四个字段，对两位导演完全相同：一句话的 \texttt{stage\_objective}（阶段目标）；\texttt{why\_now}（为何是现在），以导演在游玩中所见为依据；一到三条 \texttt{priorities}（优先事项），以方向而非代码修改的形式表述；以及零到三条需要 \texttt{preserve}（保留）的内容。循环在三处使用简报——规划器、阶段排序，以及改进和美术提示词——并且监控器据以打分的阶段清单会被重建，使简报自身的准则占据最大权重（附录~\ref{app:monitor}）。

\begin{figure}[t]
\centering
\fig{9}
\caption{简报表单，即图~\ref{fig:director} 页面的下半部分。模型导演通过该表单提交所用的端点，以 JSON 形式提交同样的四个字段。}
\label{fig:brief}
\end{figure}

表~\ref{tab:briefs} 将 \textit{Alley Brawlers} 的两份简报并列展示。它们针对同一棵树、相隔四分钟写成，而且毫无重叠。人类把这个构建看作呈现问题，要求改进美术和新手引导；模型打了两场完整比赛，从未让对手掉过一点血，于是把这个构建看作机制问题。二者都不是唯一正确的答案，结果也很接近：从这个检查点出发，人类引导的分支达到 62.0，模型引导的分支达到 61.5，而从同一构建重启、不给任何简报的对照分支达到 60.0。这项比较表明的是，简报提供了一个循环自身的证据在十八轮中都未曾产生的方向，并且两位导演通过同一工具看同一个构建，可以给出截然不同的方向。

\begin{table}[t]
\centering\small
\begin{tabular}{@{}p{0.11\linewidth}p{0.85\linewidth}@{}}
\toprule
\multicolumn{2}{@{}l}{\textbf{人类导演}——60/60 次操作，0 次重置。} \\
\midrule
目标 & 美术需要改进：标题画面和游戏结束画面杂乱且没有像样的背景，角色美术很丑，应该用漫画风格重新生成，攻击也没有特效。游戏还需要教会玩家怎么玩——现在我根本不知道该按什么。 \\
为何是现在 & 角色能动了，也有了标题和游戏结束画面，但它仍然太粗糙，缺乏吸引力。 \\
\midrule
\multicolumn{2}{@{}l}{\textbf{模型导演（Claude Opus 5）}——25/60 次操作，0 次重置。} \\
\midrule
目标 & 把比赛从一场按剧本进行的单方面痛殴变成真正的双方对战：玩家的输入必须可见且可靠地命中、格挡和脱身，对手是一个可以被击败的 AI，两名格斗者在外观和打法上确实不同。 \\
为何是现在 & 在两场完整的比赛中（两个格斗者各用一次），AI 一个像素的血都没掉。每个回合都在大约 18 秒内结束，我被压在左边角落里陷入“命中、倒地、命中”的循环。轻攻击、重攻击、特殊指令输入、按住后退、下蹲和跳跃，全都无法换血、格挡或让我脱身。当按什么都不起作用时，没有人能学会玩这个游戏。[……] 七轮视觉润色都花在了打斗周围的框架上（标题、背景、横幅、结算画面），但打斗本身仍然缺席。 \\
\bottomrule
\end{tabular}
\captionsetup{position=bottom}
\caption{为同一个 \textit{Alley Brawlers} 构建（树 \texttt{c12fe90b30b6c176}）写下的两份简报，节选其前两个字段。限于篇幅，省略了优先事项和保留项。}
\label{tab:briefs}
\end{table}

\begin{figure}[t]
\centering
\fig{10}
\caption{导演打开的案例列表：六个游戏、各自正在评审的构建，以及剩余预算。案例标识符为内部标签。}
\label{fig:caselist}
\end{figure}

\subsubsection{案例集}
共有六个游戏接受了评审，每个游戏都在其自身运行收敛到的构建上接受评审，范围从刚生成的项目到第十八轮不等（图~\ref{fig:caselist}）。每个游戏都由人类导演和模型导演各自独立评审，随后每份简报都交给一个从同一构建重启的分支，并与一个不给简报的对照分支进行比较。六个游戏中有三个是附录~\ref{app:cases} 的案例研究。

\subsection{智能体提示词与输出约定}\label{app:prompts}
本附录报告承担循环决策的五个模块的提示词，每个模块对应第~3 节的一个阶段。每份提示词都被删节为其决策骨架：固定的目录、示例、画面标注约定和框架内务性内容都被删去，每处删减都标记为 \texttt{[...]}。诸如 \texttt{\{items\}} 的花括号是运行时插槽，由开发状态、试玩记录或回放画面填充。

有一个设计选择在全部五个模块中反复出现，统一说明一次更容易看清。每份提示词都被告知回答其自身问题所需的信息，而刻意不被告知其余信息。规划器得到的本轮方向是一个结论，而不是一个论证。验证器被告知本轮试图做什么，因为那正是它必须检验的主张；而比较器既不被告知有任何东西被改进过，也不被告知哪个构建更新。编辑器被告知问题已经确立，这样它就从“这在代码的哪里”开始，而不是重新推导诊断。正是这种“有所保留”，使每个阶段不会去确认自己的预期。

\begin{table}[t]
\caption{\rsig{} 循环中五个面向 LLM 的模块，按其实现的第~3 节阶段排列。}
\label{tab:modules}
\centering\footnotesize
\begin{tabular}{@{}p{0.15\linewidth}p{0.11\linewidth}p{0.36\linewidth}p{0.3\linewidth}@{}}
\toprule
阶段 & 模块 & 输出约定 & 使用方 \\
\midrule
方向决策 & 控制器 & 一个 \texttt{development\_question}；动作取自 \{\texttt{reuse}, \texttt{explore}, \texttt{escalate}\}；一份 \texttt{exploration\_plan}，指明能回答该问题的最低成本方式。 & 设定本轮目标；\texttt{escalate} 报告该方向已耗尽，而不是凭空编造工作。 \\
智能体探索 & 探索器 & 一次会话应解决的 1–3 条需求、要启动的场景，以及给游玩智能体的一条可执行指令。 & 驱动测试智能体与 $P_t$ 交互，产生轨迹 $\tau_t$。 \\
基于证据的编辑 & 编辑器 & 在构建约束下对项目的编辑，以及一份修改内容报告。 & 产生 $P_{t+1}$。 \\
智能体验证 & 验证器 & \texttt{goal\_achieved} 取自 \{\texttt{yes}, \texttt{no}, \texttt{unclear}\}；\texttt{broken\_now}；当该列表非空时给出 \texttt{next\_goal}。 & 接受或拒绝本轮；任何被破坏的内容都作为待办工作重新进入状态。 \\
全局质量监控 & 全局质量监控器 & 两个构建的逐准则读数、带方向和幅度的变化，以及退化情况——从不给出分数。 & 汇总为决定冠军的代理值（附录~\ref{app:monitor}）。 \\
\bottomrule
\end{tabular}
\end{table}

\subsubsection{控制器}
规划器决定的是要\emph{查明}什么，而从不决定要\emph{改变}什么。它得到的输入包括：作为既定结论的本轮方向、需求状态、构建上次被游玩时的表现、可用的演示以及剩余预算；它返回一个问题以及回答该问题的最低成本方式。提示词的大部分篇幅用于排除那些任何观察都无法解决的问题。

\begin{promptbox}{pbg1}{提示词 P1. 控制器}
\begin{PV}
你负责规划一个正在逐轮开发的游戏接下来要调查什么。

你不写代码、不提出补丁，也不说应该改什么。那是另一个
智能体的工作，在你要求的证据收集完之后进行。你的输出
是关于要“查明”什么的决定。

[... 方向、需求状态、上次试玩、演示、预算 ...]

====================
要决定的内容
====================
一个主要问题。它必须：
  - 位于上述方向之内；
  - 以上述状态和证据为依据，而不是以这类游戏通常
    具备什么为依据；
  - 能通过观看游戏来回答：有人回放它并观察。而不是
    通过阅读源代码——没有人会替你去读；
  - 足够具体，使你能说出那个观察结果会是什么；
  - 有助于判断某处是否需要修复。

不是问题：“游戏能被改进吗？”——没有任何观察能回答它。
不是问题：“改进 HUD 布局。”——那是一条修改指令，而且
目前还什么都没有观察到。

然后选择一个动作：
  reuse      上面的证据已经说明了问题是什么。
  explore    证据缺失或含糊，而某个具体观察能解决它。
  escalate   这个方向上已没有有依据的未决问题。
             如实说明，而不是编造一个。

如果选择探索，选出能回答该问题的最低成本方式：
  replay_existing_demo    已有演示恰好覆盖了你需要看的内容
  extend_existing_demo    某个演示能到达该状态，但在得出
                          答案前就停止了
  custom_targeted_probe   没有演示能到达。代价最高的选项。

如果某个问题已经连续问了好几轮仍未解决，用同样的方式
再问一次就不算计划。要么说明这次会有什么不同，要么
换一个问题。

只用严格的 JSON 回复，不要输出其他内容：
{"action": "reuse|explore|escalate",
 "development_question": {"question": "...", "why_now": "..."},
 "exploration_plan": {"mode": "...", "demo_id": "...", "focus": "...",
                      "stop_when": "..."}}
\end{PV}
\end{promptbox}

\subsubsection{探索器}
探索是有目标的，而非泛泛而为。探针选择器把本阶段尚无法回答的需求转化为一条游玩智能体可以执行的指令，并被要求在重要性与一次短会话能到达的范围之间进行权衡。

\begin{promptbox}{pbg2}{提示词 P2. 探索器}
\begin{PV}
你要决定下一次“单次”游玩会话应该去查明这个游戏的
什么。

以下是本阶段仍无法回答的需求，以及每条可以在哪里检查：

[... 候选需求，附状态和场景 ...]

选出“此刻”最重要的 1 到 {cap} 条，并为将要游玩该构建
的人写一条指令。判断重要性——游戏其余部分所依赖的需求
胜过细节，几步操作就能解决的胜过需要长时间刷的。

如果你选的几条可以在同一个地方观察，请说明；一次需要
启动两个不同状态的会话，会把一半步数花在路上。

指令必须“可执行”：说出要做的输入或要到达的画面，以及
要观察什么。“检查飞机是否会转向”太弱。“在一次出击中，
先按左再按右各约一秒，报告飞机的航向是否改变”才是可以
执行的。

只用 JSON 回复，不要输出其他内容：
{"items": ["T04"], "scenario": "", "ask": "...", "why": "一句话"}
\end{PV}
\end{promptbox}

\subsubsection{编辑器}
编辑器接收的是一个已经确立的问题，提示词的首要任务是阻止它再去确立一遍。它获得源代码、必须通过的构建命令，以及运行并观察构建的方法；它没有修改框架或评估轨迹的权限。

\begin{promptbox}{pbg3}{提示词 P3. 编辑器}
\begin{PV}
你在改进一个小游戏。这是第 {r} 轮，最多 {R} 轮。

你没有玩过这个构建。是别人玩的：一个游玩智能体驱动了
它，画面已被解读，下面是这些证据所支持的内容以及为本轮
选定的内容。

所以下面的问题已经确立。你不需要查明它是否真实，也不
应该重新回放游戏去确认——那项工作已经完成，重复它会
耗掉你修复所需的调用次数。你的工作从“这在代码的哪里”
开始。

[... 本轮的证据包、近期历史、项目地图 ...]

你拥有的
  - 源代码，就在你工作的目录中
  - `{build_cmd}` —— 你停止时它必须通过
  - 你可以运行构建并观察它：[...]

规则
  - 改游戏，不要改框架。不要编辑 `node_modules`、`dist`、
    `demo_outputs/`，或 `_repair_evidence/` 下的任何内容。
  - 不要编辑 `demo_outputs/`。那是评估器回放的输入轨迹；
    改动它们改变的是考卷，而不是游戏。
  [... 流程与服务器规范 ...]
\end{PV}
\end{promptbox}

\subsubsection{验证器}
编辑之后，每个随附的演示都会被回放，由一个阅读者判定本轮是否配得上提交。提示词区分了两种容易混淆的失败：游戏从未有过的功能不是缺陷；一个产生了可见效果的输入已经被接收，即使脚本始终没有到达它原本设定的结果。只有毫无改变的输入才算被破坏。

\begin{promptbox}{pbg4}{提示词 P4. 验证器}
\begin{PV}
一个修复智能体刚刚编辑了一个游戏。你的工作是看看游戏
现在的表现并回答两个问题。你是这次编辑唯一的检查：
在你之后，它就发布了。

本轮试图做什么
{goal}

[... 该目标所源自的需求；回放画面 ...]

按以下顺序回答这两个问题。

1. goal_achieved —— yes / no / unclear。
   本轮打算做的事情现在是否明显成立？“unclear”是一个
   真实的答案：当演示从未到达该目标涉及的那部分游戏时，
   就这样回答，而不要猜测。

2. broken_now —— 一个列表，可以为空。
   观看这些画面的玩家会明显觉得不对的任何表现。具体说出
   你“看到”了什么、在哪一帧。最重要的几类，因为它们在
   源代码里看不出来：
     - 输入完全没有引起变化——前后画面一次又一次相同，
       说明游戏没有在接收输入
     - 有东西盖住了游戏区域：某个精灵、特效或面板挡住了
       角色、HUD 或动作
     - 画面冻结：最后一帧与第一帧相同
     - 文字无法阅读、被截断或相互重叠
   不要在这里列出缺失的功能。游戏从未有过的功能属于任务
   需求，而不是被破坏的东西。

还有一件不要报告的事。一个明显“产生了”效果的输入——
出现了模式标签、显示了提示——已经被接收，即使演示始终
没有到达它原本设定的结果。那是游戏需要的步骤比这个固定
脚本执行的更多，并不是游戏的错误。只有“毫无”改变的
输入才属于 broken_now。

然后，仅当 broken_now 非空时，写下 next_goal：用一句话
告诉修复智能体首先修什么。修最严重的那个，而不是全部。
\end{PV}
\end{promptbox}

\subsubsection{全局质量监控器}
比较器是整个方法中唯一读取像素的组件。它看到的是在同一固定脚本下回放的冠军和候选，并被问到发生了什么变化；它既不被告知哪个构建更新，也不被告知发生过改进。提示词的大部分内容用于防范成对比较出错的两种方式：假定更新的构建更好，以及把录制时序误认为行为变化。

\begin{promptbox}{pbg5}{提示词 P5. 全局质量监控器}
\begin{PV}
同一个游戏的两个构建，用“相同”的固定输入脚本回放。
你的工作是说出这个场景中发生了什么变化，别无其他。

[... 场景、它能涉及的准则、配对的画面 ...]

脚本是固定的，所以两个构建在相同时刻收到相同的输入。
对于下面的每个时刻，你会先看到旧构建在输入前后的画面，
再看到新构建在同样两个时刻的画面。

录制并非逐帧精确：同一构建回放两次，截取同一时刻也可能
相差一两帧，仅这一点就会让分数读数差一分、动画看起来
慢半拍，或让下落的方块低一行。那是录制器，而不是构建。

正是出于这个原因，每次输入后你会得到两帧。到第二帧就
消失的差异，或被之后的时刻抵消的差异，属于时序。在下一
个时刻和“最后”一帧仍然存在的差异，才属于构建。对第一
类请回答 `same`。

规则
- 只比较这些画面所支持的内容。
- 不要假定更新的构建更好，也不要假定一个你无法解释的
  差异就是缺陷。两种倾向都是错误。
- 如果某个行为在这里没有被触发，它的准则就是
  `unobserved`。没看到某样东西，不等于看到它缺失。
- 任何变差的地方都要单独报告，无论其余部分多么微小。
- 判断相对变化。不要给任何一个构建打分。
- 不要提出代码编辑、修复方案或下一步。那不是你的工作。
\end{PV}
\end{promptbox}

\section{实验协议与评估可靠性}\label{app:exp}
\subsection{实验配置}\label{app:config}
本附录在一处列出：哪个模型扮演哪个角色、各自如何部署和解码，以及一次运行所获得的全部预算。表~\ref{tab:models} 涵盖模型，表~\ref{tab:params} 涵盖循环和评分参数。

\paragraph{评分协议。}表~\ref{tab:main} 中的每个数字都是 140 个任务上的平均值，每个任务的分数由 Qwen3.8-27B 评审对同一产物进行三次独立的回放加打分运行得到。生成器未产出项目的任务记 0 分；开发尚未完成的任务也记 0 分。基线行对冻结的 $P_0$ 打分，Play2Code 行对三十轮所交付的版本打分，\rsig{} 行则对饱和停止触发时全局质量监控器所持有的版本打分。

\begin{table}[htbp]
\caption{\textbf{按角色划分的模型。}生成器因行组而异；其余每个角色在每一行中都是同一个模型，因此行与行之间的比较是开发方法的比较，而不是其中所用模型的比较。凡是答案为关于游戏的论断（坐标、判定、评分细则条目）时，解码均采用贪心方式。部署方式对成本和可复现性都有影响——OpenRouter 上的同一个模型由十几家提供商出售，因此评审和监控器固定了各自的提供商。}
\label{tab:models}
\centering\small
\begin{tabular}{@{}llll@{}}
\toprule
角色 & 模型 & 部署方式 & 解码 \\
\midrule
生成器 & \makecell[l]{表~\ref{tab:main} 中每个\\行组各一个} & \makecell[l]{框架\\默认} & \makecell[l]{框架\\默认} \\
控制器 & GLM-5.3-Flash & OpenRouter & $T=0$，top-$p$ 1 \\
探索器 & Qwen3.8-27B & 本地 vLLM（TP $=2$） & $T=0$，top-$p$ 1 \\
验证器 & GLM-5.3-Flash & OpenRouter & $T=0$，top-$p$ 1 \\
编辑器 & GLM-5.3-Flash & OpenRouter & $T=0$，top-$p$ 1 \\
全局质量监控器 & Qwen3.8-Flash & OpenRouter & $T=0$ \\
素材生成 & Seedream-5.0-lite & OpenRouter（图像） & — \\
评估评审 & Qwen3.8-27B & 本地 vLLM & \makecell[l]{$T=0$，关闭思考，\\输出 $\leq$2048} \\
\bottomrule
\end{tabular}
\end{table}

\subsection{Token 与成本核算}\label{app:cost}
\paragraph{范围。}我们对每个任务统计两个阶段。\emph{生成}是生成器产生冻结基础 $P_0$ 的单次智能体运行。\emph{开发}是开发轮次中的每一次模型调用：决定方向的控制器、游玩游戏的探索器、检查游戏的验证器、编辑游戏的编辑器，以及全局监控器的周期性检查。所有方法都不计入评估评审的打分。

\paragraph{Token。}Tok. 一列统计计费 token：不含缓存读取的输入，加上输出。缓存读取不计入这一统计，因为长时间的智能体会话在每次调用时都会重读同一前缀；但它们仍然按缓存价格计入 Cost。

\paragraph{成本。}每次调用，
\[
\text{cost} = n_{\text{in}}\,p_{\text{in}} + n_{\text{cache}}\,p_{\text{cache}} + n_{\text{out}}\,p_{\text{out}},
\]
其中 $n_{\text{in}}$ 为未缓存的输入 token，价格采用 OpenRouter 的标价（表~\ref{tab:prices}）。开发阶段并不按此方式定价，而是汇总 API 随每次调用返回的金额——代理会为控制器、探索器、验证器、编辑器和监控器的每次调用记录这一金额；上述公式用于只保留了 token 计数的场合，即生成阶段。Cost 只涵盖开发；生成成本在表~\ref{tab:gen} 中单独报告。生成成本主要由框架而非模型决定：同一个 GPT-5.5 在 Codex 下的缓存命中率为 94\%，在 OpenGame 下为 42\%，由此产生的账单相差八倍。

\begin{table}[htbp]
\caption{\textbf{循环、监控器与评分参数。}每个参数取一个值，在本文报告的每次运行中都相同。}
\label{tab:params}
\centering\small
\begin{tabular}{@{}llr@{}}
\toprule
& 参数 & 取值 \\
\midrule
\multirow{5}{*}{轮次} & 编辑器每轮可进行的工具调用次数 & 26 \\
& 每个改进轮次的墙钟时间预算 & 1800 s \\
& 每轮放弃前的改进尝试次数 & 3 \\
& 每个声称的修复的验证采样次数 & 3 \\
& 提出编辑前的探针步数 & 8 \\
\midrule
\multirow{3}{*}{运行} & 每次运行生成的素材数 & 20 \\
& 保存的项目快照 & 每轮 \\
& 保存的回放录像 & 每 3 轮 \\
\midrule
\multirow{5}{*}{监控器} & 检查点间隔 & 每 3 轮 \\
& 每个检查点回放的演示数 & $\leq$3 \\
& 比较时每个演示读取的帧数 & 16 \\
& 候选必须超过冠军的幅度 & 0.05 \\
& 饱和停止前冠军保持不变的检查点数（$K$） & 3 \\
\midrule
\multirow{3}{*}{评分} & 每个产物的回放加打分次数 & 3 \\
& 每个演示均匀采样的帧数 & 16 \\
& 每个任务的细则条目数（$M/D/V/A$） & 12–21（中位数 18） \\
\bottomrule
\end{tabular}
\end{table}

\paragraph{采用哪个价格。}OpenRouter 上的一个模型由十几家提供商提供服务，价格相差可达两倍，而模型页面只引用其中一家的报价，对若干模型而言那恰好是全场最贵的。因此，我们按每个模型最便宜的标准端点定价（表~\ref{tab:prices}）。有两项排除使这一做法定义明确。延迟分级并不是同一服务的不同卖家，因此不予考虑：flex 分级通过延后请求使 GPT-5.5 价格减半。并且，各端点是按这些运行在其上实际会花费多少来比较的，而不仅仅比较输入价格，因为报价中缓存价格为零意味着该端点不提供提示缓存，而不是缓存读取免费——若按字面理解，一家不提供缓存的提供商就会成为一个缓存占比 42–96\% 的工作负载的最便宜卖家。

\paragraph{数据来源。}生成阶段的用量来自每次生成运行自身的记录（输入、缓存读取和输出 token），随基础游戏语料一同发布；对于 Godot 上的 Codex + GPT-5.5，只记录了语料层面的均值。开发阶段的用量由智能体与 API 之间的代理逐次调用记录。

\begin{table}[htbp]
\caption{\textbf{OpenRouter 标价（美元 / 百万 token），获取于 2026-09-16。}每个模型按其在本文运行的缓存比例下最便宜的标准端点定价；之所以注明提供商，是因为同一模型在不同提供商之间的价格相差可达两倍。}
\label{tab:prices}
\centering\small
\begin{tabular}{@{}llrrr@{}}
\toprule
模型 & 提供商 & 输入 & 缓存读取 & 输出 \\
\midrule
GPT-5.5 & OpenAI & 5.000 & 0.500 & 30.00 \\
Kimi-K2.6 & Baidu & 0.408 & 0.069 & 1.72 \\
Qwen3.8-27B & DeepInfra & 0.150 & 0.037 & 1.88 \\
GLM-5.3-Flash & DeepInfra & 0.075 & 0.015 & 0.25 \\
\bottomrule
\end{tabular}
\end{table}

\begin{table}[htbp]
\caption{\textbf{冻结基础 $P_0$ 的生成，每个任务的均值。}“计费”即基线行的 Tok.；“总计”另加缓存读取；“成本”将表~\ref{tab:prices} 应用于全部三种 token 计数。两行 GPT-5.5 是同一模型在两种框架下的结果，价格不同：Codex 会重读缓存的前缀，OpenGame 则不会。统计范围为全部 140 个任务。}
\label{tab:gen}
\centering\small
\begin{tabular}{@{}llrrrrr@{}}
\toprule
引擎 & 生成器 & 计费 & 总计 & 缓存 & 输出 & 成本 \\
\midrule
Godot & Codex + GPT-5.5 (high) & 0.26M & 3.83M & 94\% & 28k & \$\,3.77 \\
Godot & GLM-5.3-Flash & 0.44M & 11.43M & 96\% & 36k & \$\,0.20 \\
Godot & Kimi-K2.6 & 3.22M & 8.98M & 65\% & 123k & \$\,1.87 \\
Godot & Qwen3.8-27B & 6.41M & 11.62M & 46\% & 179k & \$\,1.47 \\
Godot & Qwen3.8-27B (SFT) & 0.57M & 8.21M & 95\% & 140k & – \\
Phaser & OpenGame + GPT-5.5 & 5.36M & 9.21M & 42\% & 97k & \$31.17 \\
Phaser & Qwen3.8-27B & 9.03M & 26.69M & 67\% & 292k & \$\,2.51 \\
\bottomrule
\end{tabular}
\end{table}

\subsection{回放与评分稳定性}\label{app:stability}
\paragraph{回放变异性。}GameCraft-Bench 依据观察到的游戏过程而不仅仅是源代码来为游戏打分。即使对于同一个冻结项目，由于游戏执行、输入调度和画面捕获并不完全同步，重复回放也可能产生略有不同的轨迹。这种变异的幅度取决于游戏，表~\ref{tab:family-var} 测量了它取决于哪类游戏。因此，对同一段录像反复评审并不能刻画评估过程的全部不确定性。

\paragraph{评估协议。}为了应对这种变异性，每个产物都经过三次独立的端到端运行评估。每次运行都从一次全新的回放开始，生成自己的游戏录像，并由本文所有分数都使用的同一个 Qwen3.8-27B 评审独立打分。我们报告三次运行的平均分。评审配置在所有评估中保持固定，因此重复运行主要刻画的是游戏回放引入的变异，而不是评估设置的变化。

\paragraph{观察到的稳定性。}表~\ref{tab:stability} 汇总了 529 个冻结产物上的变异，即本文中每个被打分三次的 Godot 产物：两个生成器的冻结基础、Play2Code 交付的版本，以及 \rsig{} 交付的版本。

大多数产物接近确定性。三次独立运行间极差的中位数为 0.73 分，529 个中有 225 个三次得分完全相同，四分之一的产物变化不到百分之一分。改进后的产物比冻结的产物变化略大（中位数 0.87 对 0.36），这正是一个包含更多游戏内容的构建应有的表现。

\begin{table}[htbp]
\caption{\textbf{同一冻结产物三次独立端到端评估的得分变异。}每次运行使用全新的游戏回放，并由 Qwen3.8-27B 评审重新评判。极差指三次运行中最高与最低总分之差；各列为其四分位数。}
\label{tab:stability}
\centering\small
\begin{tabular}{@{}lrrrrr@{}}
\toprule
子集 & $n$ & p25 & 中位数 & p75 & 三次相同 \\
\midrule
全部产物 & 529 & 0.00 & 0.73 & 3.33 & 225 \\
基础 $P_0$ & 266 & 0.00 & 0.36 & 2.33 & 131 \\
改进后 & 263 & 0.00 & 0.87 & 3.50 & 94 \\
\bottomrule
\end{tabular}
\end{table}

\begin{table}[htbp]
\caption{\textbf{按任务类别的回放变异}，仅列出至少有十二个已打分产物的类别，按极差中位数排序。产物集合与极差定义同表~\ref{tab:stability}。}
\label{tab:family-var}
\centering\small
\begin{tabular}{@{}lrr@{\hspace{3em}}lrr@{}}
\toprule
类别 & $n$ & 中位数 & 类别 & $n$ & 中位数 \\
\midrule
sports & 14 & 4.66 & shooter & 24 & 0.44 \\
horror & 19 & 3.37 & puzzle & 32 & 0.42 \\
openworld & 62 & 1.41 & simulation & 24 & 0.28 \\
visualnovel & 32 & 1.17 & strategy & 63 & 0.00 \\
platformer & 77 & 0.87 & rhythm & 21 & 0.00 \\
tycoon & 58 & 0.73 & idle & 16 & 0.00 \\
roguelike & 53 & 0.62 & racing & 16 & 0.00 \\
cardgame & 18 & 0.50 & & & \\
\bottomrule
\end{tabular}
\end{table}

\paragraph{哪些游戏会变化。}按任务类别拆分这 529 个产物，可以定位变异的来源，而不是把它笼统地归为“游戏动态”的属性（表~\ref{tab:family-var}）。划分并非按快慢：变化最大的三个类别是 sports、horror 和 openworld，而 racing 和 rhythm——它们与基准中任何游戏一样都是实时的——极差中位数恰好为零，strategy 亦然。变化大的类别的共同点在于，它们的评分细则答案出现得较晚。一场比赛是否获胜、一次闹鬼事件是否解决、一个世界是否被穿越，取决于二十秒的游玩恰好结束在哪里；而一个节拍是否准时落下、一个圈速计时器是否运行，则取决于回放每次都能重现的那些帧。变异追随的是某条准则的证据在会话中出现得有多晚，而不是游戏运动得有多快。

这也说明了单次运行比较在哪里最不可靠。一分的差异对节奏游戏是有意义的，对体育游戏则是噪声，这也是阅读附录~\ref{app:family} 的按类别表格时应结合该类别变异的原因之一。

\paragraph{跨评审稳健性。}为了检验主要结果是否依赖于某个特定的评估模型，我们用 GPT-5.5 重新为 40 个按类别分层抽样的 Godot 任务打分，使用与 Qwen3.8-27B 评审完全相同的回放、采样帧、任务特定细则和评分公式；只更换评审模型。

如表~\ref{tab:crossjudge} 所示，两个评审给出相同的方法层面排序：\rsig{} $>$ Play2Code $>$ 基线，且改进幅度相近。\rsig{} 相对冻结基线的增益在 Qwen3.8-27B 下为 +14.04 分，在 GPT-5.5 下为 +12.04 分；相对 Play2Code 的增益分别为 +13.47 和 +11.76；四个 bootstrap 区间都远高于零。在 40 个任务中，两位评审对“\rsig{} 优于基线”这一比较的方向在 37 个任务上一致。尽管它们的绝对分数仅中等相关（$\rho = 0.64$），但从相对比较中得出的结论在不同模型家族间保持稳定。

\begin{table}[htbp]
\caption{\textbf{40 个按类别分层的 Godot 任务上的跨评审稳健性。}Qwen3.8-27B 和 GPT-5.5 使用相同的细则为相同的回放录像打分。两位评审都保持了方法排序，并给出相近的成对改进幅度。置信区间为 95\% 任务级 bootstrap 区间。}
\label{tab:crossjudge}
\centering\small
\begin{tabular}{@{}lcc@{}}
\toprule
& \multicolumn{2}{c}{总分} \\
\cmidrule(l){2-3}
方法 & Qwen3.8-27B & GPT-5.5 \\
\midrule
基线（冻结的 $P_0$） & 49.07 & 50.34 \\
\quad + Play2Code & 49.63 & 50.62 \\
\quad + \rsig & \textbf{63.10} & \textbf{62.38} \\
\bottomrule
\end{tabular}

\vspace{6pt}
\setlength{\tabcolsep}{4pt}
\begin{tabular}{@{}lrrccc@{}}
\toprule
比较 & $\Delta$ Qwen & $\Delta$ GPT-5.5 & 95\% CI Qwen & 95\% CI GPT-5.5 & 一致数 \\
\midrule
\rsig{} $-$ 基线 & +14.04 & +12.04 & [+9.84, +18.59] & [+9.53, +14.60] & 37/40 \\
\rsig{} $-$ Play2Code & +13.47 & +11.76 & [+8.95, +18.19] & [+7.79, +16.26] & 34/40 \\
Play2Code $-$ 基线 & +0.56 & +0.28 & [$-$3.45, +4.21] & [$-$3.83, +3.83] & 24/40 \\
\bottomrule
\end{tabular}

\vspace{2pt}
{\footnotesize 一致数统计两位评审偏好同一方法的任务数。}
\end{table}

\subsection{自由游玩评估}\label{app:freeplay}
\paragraph{协议。}GameCraft-Bench 通过回放随提交项目打包的演示轨迹来评估游戏。这些轨迹是可见的开发资源，而不是隐藏的基准轨迹，\rsig{} 会复用它们在检查点之间进行受控的前后比较。尽管隐藏的评分细则、分数和评审反馈在开发过程中从不暴露，反复使用相同的演示仍可能偏向针对这些交互轨迹的特定改进。因此，我们评估这些收益在独立生成的游戏过程中是否依然存在。

我们抽取 20 个按类别分层的 Godot 任务。对每个任务，冻结基线、Play2Code 和 \rsig{} 的构建分别由盲测智能体独立游玩，每个智能体有 40 次操作和一次重置。智能体只能看到公开的游戏规格说明和正在运行的游戏——看不到提交的演示、评分细则、分数、源代码、开发历史或方法身份——并返回关于其观察的事实性报告。另一个评审对这些报告进行成对比较，共得到 60 次比较。

\paragraph{评估。}成对评审遵循主评估所用的四个质量维度和相对权重，同时只依据自由交互中观察到的证据进行评判。重要的是，回放脚本从不向游玩智能体展示：它们自行选择动作、时机和轨迹。这改变了交互分布，同时保持游戏质量的定义与主评估一致。

\paragraph{结果。}如表~\ref{tab:freeplay} 所示，\rsig{} 在 20 个任务中的 17 个上优于冻结基线、在 18 个上优于 Play2Code，而 Play2Code 相对基线并未显示出明显优势。在所有有胜负的比较中，自由游玩评估与基准排序在 59 例中的 45 例上一致。这些结果表明，\rsig{} 的收益在独立选择的游戏轨迹下依然存在，而不局限于开发和基准评估中使用的固定回放脚本。

\begin{table}[htbp]
\caption{\textbf{20 个按类别分层的 Godot 任务上的自由游玩评估。}盲测智能体独立游玩每个构建，另一个评审使用与主评估相同的质量维度和权重比较它们的报告。$p$ 为对有胜负比较所做的双侧符号检验。}
\label{tab:freeplay}
\centering\small
\begin{tabular}{@{}lccc@{}}
\toprule
比较 & 胜/负/平 & $p$ & 一致数 \\
\midrule
\rsig{} – 基线 & 17/3/0 & 0.003 & \multirow{3}{*}{45/59} \\
\rsig{} – Play2Code & 18/2/0 & $<$ 0.001 & \\
Play2Code – 基线 & 9/11/0 & 0.824 & \\
\bottomrule
\end{tabular}

\vspace{2pt}
\parbox{0.85\linewidth}{\footnotesize 胜/负/平为第一个方法的胜、负、平局数。一致数统计自由游玩评估与基准得分偏好同一构建的有胜负比较数。}
\end{table}

\subsection{主要结果的统计可靠性}\label{app:stat}
\paragraph{配对的任务级不确定性。}主表中的所有比较都按任务配对。我们对 140 个基准任务进行 20,000 次 bootstrap 重采样，以估计总分平均差 $\Delta$ 的不确定性，并额外报告对配对任务级差异的 Wilcoxon 符号秩检验。我们还单独对每个固定任务的三次回放加打分运行进行重采样，以衡量评估噪声；这些区间远窄于任务 bootstrap 区间，说明基准任务之间的差异主导了回放和评审噪声。

\paragraph{主要比较。}在表~\ref{tab:sig} 包含的每种引擎–生成器设置下，\rsig{} 都显著优于其冻结的起点，增益为 +8.8 到 +14.3 个总分点，任务 bootstrap 区间远高于零。与 Play2Code 相比，改进在两个 Godot 生成器以及 Phaser 上的 GPT-5.5 中都是显著的；Phaser 上 Qwen3.8-27B 的比较仍不能确定（+1.55，$p = 0.394$）。

\begin{table}[htbp]
\caption{\textbf{主要结果的配对显著性。}$\Delta$ 为任务级总分的平均差。置信区间来自对基准任务进行的 20,000 次配对 bootstrap 重采样；$p$ 来自 Wilcoxon 符号秩检验。胜/负统计第一个方法得分更高/更低的任务数。}
\label{tab:sig}
\centering\small
\begin{tabular}{@{}lrccc@{}}
\toprule
比较 & $\Delta$ & 95\% CI & $p$ & 胜/负 \\
\midrule
\multicolumn{5}{@{}l}{\textit{Godot}} \\
GPT-5.5：\rsig{} – 基线 & +14.26 & [11.49, 17.06] & $<10^{-4}$ & 112/16 \\
GPT-5.5：\rsig{} – Play2Code & +13.79 & [10.89, 16.65] & $<10^{-4}$ & 118/19 \\
Qwen3.8-27B：\rsig{} – 基线 & +10.70 & [7.75, 13.77] & $<10^{-4}$ & 100/31 \\
Qwen3.8-27B：\rsig{} – Play2Code & +7.24 & [4.49, 10.04] & $<10^{-4}$ & 92/41 \\
\midrule
\multicolumn{5}{@{}l}{\textit{Phaser}} \\
GPT-5.5：\rsig{} – 基线 & +8.77 & [6.13, 11.51] & $<10^{-4}$ & 101/35 \\
GPT-5.5：\rsig{} – Play2Code & +3.11 & [0.53, 5.61] & 0.008 & 83/53 \\
Qwen3.8-27B：\rsig{} – 基线 & +10.21 & [6.99, 13.70] & $<10^{-4}$ & 95/36 \\
Qwen3.8-27B：\rsig{} – Play2Code & +1.55 & [$-$1.50, 4.66] & 0.394 & 73/59 \\
\bottomrule
\end{tabular}
\end{table}

\subsection{自适应方向策略评估}\label{app:adaptive}
\paragraph{实验设置。}我们在 16 对从相同冻结基础项目初始化的游戏上，比较自适应开发与轮转方向策略。每种策略运行 30 个开发轮次，使用相同的智能体模型、工具和每轮预算；唯一的区别在于如何选择开发重点。图~\ref{fig:adaptive}a,b 使用在开发过程中对两种策略都不可见的信号来分析所得轨迹。我们将“构建失败”定义为所得项目无法产出可运行游戏的轮次；随后的轮次相对于该失败事件对齐到偏移 +1。对于按状态条件的分析，我们另外使用一份保留的诊断细则来识别当前落后的质量维度。该细则仅用于事后分析，从不提供给开发循环。图~\ref{fig:adaptive}a 中的阴影区域为相应失败事件或分组开发轮次上的 $\pm1$ 标准误。

\paragraph{盲测成对评估。}为比较两种策略产出的最终游戏，我们使用它们第 30 轮的检查点进行了 63 次盲测成对比较。每次比较中，同一条演示轨迹在两个构建上回放，所得观察并排呈现给评估者。左右位置在每次比较中独立随机分配，且不透露任何关于底层方向策略的信息。评估者报告总体偏好，以及在机制、深度、视觉和美术各维度上的偏好；允许平局。为检查位置偏差，一部分比较在交换左右后被再次评估。这些重复判断仅用作一致性检查，不计入报告的 63 次比较。成对评估者独立于 GameCraft-Bench 基准评估器，无法访问其分数、细则或反馈。

\subsubsection{成对评判准则}
以下完整给出交给第~5 节盲测成对评审的准则。它们刻意不同于基准的评分细则：评审依据画面所示比较两个构建，而基准则依据任务为单个构建打分。规则 1 起的作用最大——没有它，一个坏掉的构建也能凭它确实渲染出来的那些帧拿到分数。

\begin{promptbox}{pbg5}{比较细则.\ 盲测成对评审}
\begin{PV}
你是一名资深游戏评测者。你要比较同一个游戏的两个构建，
游戏 1 和游戏 2。二者都由同一个起始游戏和同一份设计文档
开发而来，设计文档见下文。每个构建都用完全相同的脚本化
输入（一个“演示”）游玩过，你看到的是从该录像中以相同
固定间隔、按时间顺序采样的画面。只依据画面中可见的内容
进行评判。

对下面的每个维度，判断哪个构建更好："1"、"2" 或 "tie"。

## 核心机制
脚本化输入是否明显产生了设计文档所描述的游戏玩法？
- 更好：操作在屏幕上有清晰的效果（移动、攻击、放置、
  状态变化、分数或资源变化）；能观察到设计文档的核心循环。
- 更差：输入没有可见效果；游戏停在标题或菜单画面；
  场景加载失败；屏幕空白或冻结。

## 内容深度
在这个演示中，设计好的游戏有多少内容出现并被触及？
- 更好：设计文档中更多不同的机制、敌人、关卡、事件、
  进程或胜负状态出现。
- 更差：单一的重复画面；占位内容；UI 中宣布了却从未
  展示的系统。

## 功能性视觉
玩家能否读懂游戏状态？
- 更好：HUD、反馈和重要物体清晰可读，并与背景明显区分；
  文字没有被截断或重叠；精灵和动画帧与帧之间稳定。
- 更差：文字无法阅读或相互重叠；精灵闪烁或不一致；关键
  物体被其他物体遮挡；重要事件没有可见反馈。

## 呈现与美术
它看起来像一个连贯、完成的游戏吗？
- 更好：各画面风格一致；构图有意为之；素材与设计文档的
  主题相符。
- 更差：该有美术的地方是基本图形；风格混杂或冲突；视觉
  噪声；让游戏更难读懂的装饰。

## 总体
综合考虑，这份设计文档的玩家会更喜欢哪个构建？

## 规则
1. 在这些画面中坏掉的构建（空白、崩溃、冻结、卡在演示
   本应离开的画面上），在它未能展示的每个维度上都判负，
   无论它其他画面看起来多好。
2. 更多素材、更多特效或更繁忙的画面本身并不更好。只有
   当它们服务于设计并保持游戏可读时才予以偏好。
3. 忽略构建显示在哪一侧。游戏 1 和游戏 2 的顺序是随机的。
4. 当两个构建在某个维度上没有明显差异，或二者同样失败时，
   回答 "tie"。
5. 每条理由都必须指向可见的内容，点明构建和帧（例如
   “游戏 2，第 7 帧：分数计数器与计时器重叠”）。
\end{PV}
\end{promptbox}

\subsection{智能体验证审计}\label{app:audit}
\paragraph{审计协议。}我们在由 GPT-5.5 初始化的 40 个游戏的开发轨迹上评估验证的可靠性。我们从这些轨迹中抽取改进前和改进后的决策，并将与每个决策相关的完整交互证据交给 Claude Opus 5\citep{anthropic2026}进行独立裁定。每个案例评估两次；两次判断之间的分歧由人工复核以得到最终标签。该评估者仅用于验证审计，与 GameCraft-Bench 评估器及其分数、细则和反馈完全隔离。

\paragraph{改进前溯据。}我们评估所提出的改进目标是否得到编辑前可用证据的支持。对每个被审计的会话，若观察到的游戏证据支持所声称的问题或改进机会，则该目标被标记为\emph{有据}；若缺乏支持或与证据矛盾，则标记为\emph{无据}。我们报告有据精确率（可判定目标中有据者的比例）以及每轮无据目标的平均数。我们将 \rsig{} 使用的智能体证据溯据流程与一个在相同开发上下文中运行的自由形式评论者进行比较。

\paragraph{改进后失败检测。}我们进一步评估在执行更新后的游戏之后，验证能否检测出不成功的编辑。当预期的变化在游戏中无法观察到，或该编辑引入了退化时，一次改进被标记为\emph{不成功}。以不成功的改进为正类，我们在 48 个经裁定的轮次上报告失败召回率、特异度和平衡准确率。我们将基于回放的智能体验证与一个仅检查构建的基线进行比较，后者只检查更新后的项目能否成功执行，而不检查其交互行为。

\subsection{多阶段指引评估}\label{app:multistage}
第~5 节检验的是开发的单个阶段。本附录说明我们如何检验第二和第三阶段是否持续带来回报：每对阶段冠军都通过盲测成对游玩进行比较，因此较晚的阶段必须作为一个游戏被偏好，而不仅仅是得分更高。

\begin{wrapfigure}{r}{0.46\linewidth}
\vspace{-10pt}
\centering
\includegraphics[page=11,width=\linewidth]{\figpdf}
\caption{每对中较晚阶段的盲测成对胜率，附 $\pm1$ 标准误。}
\label{fig:multistage}
\vspace{-8pt}
\end{wrapfigure}

\paragraph{比较什么。}三对构建，每个都是全局质量监控器在某阶段结束时交付的冠军：阶段 2 对阶段 1、阶段 3 对阶段 2，以及阶段 3 对阶段 1。一对中的两个构建来自同一个游戏，因此每次比较都在任务内部进行，双方使用相同的生成器、相同的任务规格说明和相同的每阶段预算。结果显示，较晚的阶段始终比较早的阶段更受偏好（图~\ref{fig:multistage}），表明高层指引能够在饱和之后反复重新打开开发的提升空间，并支持持续的多阶段演化。

\paragraph{为何采用成对比较。}一个修复了坏掉机制的阶段，和一个添加了无用菜单的阶段，可能使基准得分变动相同的幅度。因此，本评估问的是玩家更愿意拥有哪个构建，而不使用基准得分作为结果。

\paragraph{三个智能体，各看一个构建。}一次比较使用三个智能体。两个游玩智能体各获得一个构建，从不看到另一个，也看不到任何分数、源文件、轮次编号或开发历史；每个智能体返回一份固定格式的事实报告——游戏能否启动、核心循环能否完成、输入如何响应、渲染出了什么、什么缺失或坏掉——并被明确禁止给出评判。第三个评审智能体只阅读这两份报告（顺序在每次比较中随机化并记录），返回“报告 1 更好”“报告 2 更好”或“平局”，并附上置信等级。评审被要求将一个具体点名的缺陷看得比笼统印象更重，并将关于评审框架的任何内容——一次调用推进多少游戏时间、一个构建启动要多久、操作预算——视为两个构建共有、因而不构成对任何一方的证据。

\paragraph{预算相等。}每个游玩智能体都在 40 次操作和一次重置的固定预算下，通过附录~\ref{app:director} 的评审工具驱动其构建，其中 \texttt{wait}、\texttt{look} 和 \texttt{state} 免费。会话以评审者 id 为键，重新打开的会话会带着已花费的预算继续，因此每个智能体都获得该次比较独有的 id，并在打开后立即检查自己的操作计数器；报告计数非零的会话会被丢弃，该比较重新运行。在不相等预算下比较两个构建，得到的结论追随的是预算而不是构建，这就是为什么这项检查是协议的一部分而非事后补充。

\paragraph{报告方式。}当评审偏好属于较晚阶段的那份报告时，该比较计为较晚阶段胜出。平局单独报告，并从胜率中排除而不是对半拆分，因为平局表明两份报告无法区分这两个构建。胜率附带 Wilson 形式的 $\pm1$ 标准误，可在接近端点时使区间保持在 $[0,1]$ 之内；每对只有几十个游戏，这些区间较宽，正文也如实说明这一点，而不是从相互重叠的柱子中读出排名。

\section{训练数据与经验内化}\label{app:train}
\subsection{训练语料的构建与整理}\label{app:corpus}
发布的语料由两部分组成。\emph{基础游戏}即生成的项目本身，以及产生每个项目的运行：1,998 个去重后的游戏（1,446 个 Godot、552 个 Phaser），其中 1,070 个附有生成轨迹，另有分布在三个生成批次中的 910 棵完整的逐次试验运行树。\emph{监督数据}是我们用于训练的内容，分三个阶段从这些运行中抽取，汇总于表~\ref{tab:corpus}：S1 游戏生成、S2 规划和 S3 改进。生成轨迹由 GPT-5.5 产生，改进轨迹由 GLM-5.3-Flash 产生。规划并非自由撰写：1,237 份从已完成游戏的文件树中蒸馏得到，870 份来自更早的一个规划集，仅当其提到的每个路径都存在于对应游戏中时才予以保留。

\begin{table}[htbp]
\caption{\textbf{发布的训练语料，按已发布的数据表计数。}一个\emph{决策行}是一次助手回合连同其之前的历史，因此决策行会重复上下文，其 token 计数并不相互独立；轨迹行是同样的录制内容未经切窗的形式，应以其统计 token。经验证的改进轮次是指经独立的改进后验证判定已实现既定目标的轮次；其余轮次作为负例发布，而不是丢弃。}
\label{tab:corpus}
\centering\small
\begin{tabular}{@{}llllrrr@{}}
\toprule
阶段 & 单位 & 教师模型 & & Godot & Phaser & 合计 \\
\midrule
S1 生成 & 轨迹 & GPT-5.5 & & 1,929 & 284 & 2,213 \\
S1 生成 & 决策行 & GPT-5.5 & & 66,312 & 28,097 & 94,409 \\
S2 规划 & 规划 & 蒸馏 & & 1,979 & 129 & 2,108 \\
S3 改进 & 轮次 & GLM-5.3-Flash & & 3,730 & 273 & 4,003 \\
\quad 其中经验证 & 轮次 & — & & 1,869 & 144 & 2,013 \\
S3 改进智能体 & 决策行 & GLM-5.3-Flash & & 111,302 & 9,930 & 121,232 \\
\quad 其中经验证 & 决策行 & — & & 48,341 & 5,098 & 53,439 \\
\midrule
被改进的游戏 & 游戏 & — & & 345 & 43 & 388 \\
不同任务 & 任务 & — & & 2,189 & 322 & 2,511 \\
训练行 & 行 & — & & 183,063 & 38,689 & 221,752 \\
\bottomrule
\end{tabular}
\end{table}

\paragraph{轨迹序列化。}每个生成会话被序列化为单条多轮智能体轨迹，包含任务规格说明、助手输出、工具调用和工具结果。训练损失只施加于助手生成的输出，避免对共享的交互上下文进行重复监督。规划样本取自智能体在编码之前的初始实现计划，而之后对同一计划的进度更新被丢弃。一个决策行保留一个固定前缀（系统提示、任务，以及存在时的规划），然后在 32k token 预算内放入尽可能多的最近完整决策，最后放入目标决策。

\paragraph{过滤掉什么。}仅当会话产出了可执行产物时才保留生成轨迹；判断依据是产物是否存在，而不是智能体 SDK 报告的成功状态，因为后者并不能可靠地表明是否产出了可运行的项目。空轨迹和重复任务被移除。仅当一份规划引用的每个文件路径都存在于完成的游戏中、且其引擎与简报相符时，才保留该规划。当改进轮次的会话失败或其差异为空时，该轮次被丢弃；且仅当经过独立的改进后验证时才标记为正例——Godot 的 3,730 个轮次中有 1,869 个，Phaser 的 273 次尝试中有 144 次。冗长的工具输出（例如构建日志）会被从中间截断，保留所执行的命令及其结论。整个项目的差异被压缩为改进实际修改的源代码片段，使改进样本的中位大小缩小约 48 倍。

\paragraph{评估污染。}任何基线、消融或评估运行的内容都不会进入语料，140 个不同的基准任务 id 在构建时即被排除。发布的数据表事后又针对这些 id 重新扫描：在全部 13 种配置、共 2,511 个任务 id 的并集上，重叠为零。

\paragraph{记录缺口被标记而非隐藏。}有一个智能体 SDK 保存的是工具结果的一行 UI 摘要，而不是模型实际收到的内容。受影响的行被保留，并带有显式标记（{\small\texttt{row\_has\_summary\_only}、\texttt{target\_follows\_summary\_only}、\linebreak\texttt{tool\_results\_source}、\texttt{trajectory\_incomplete}}），因此一个干净的训练子集只是一个过滤条件，而不是另一个发布版本。这影响全部 284 条 Phaser 生成轨迹（其中 20.2\% 的行的目标决策紧跟在摘要之后），以及两个 Godot 改进分支中的一个（31.7\% 的多轮改进行）；另一个 Godot 分支是完整的。

\paragraph{数据脱敏。}每一行、每个元数据列和每份文档都逐项重写。绝对机器路径变为 \texttt{/workspace/\allowbreak game}、\hspace{0pt}\texttt{/workspace/\allowbreak harness} 或 \texttt{<path>}（403,280 处重写）；用户、主机、账户和内部项目名称变为 \texttt{<redacted>}（13,610 处）；模型网关主机、内部镜像和集群名称被替换；电子邮件地址和非回环 IP 字面量被移除；捕获的目录列表中的属主列被重写；26,577 个配置字段被就地中和，因此数据模式保持不变。转储了构建主机进程表的工具输出被直接删除而不是清洗（26 份列表，379 行）。科学内容——游戏代码和素材、细则文本、逐需求的分数、模型标识符、推理强度、耗时——均保持不变。清洗器由另一个独立编写的扫描器进行验证，后者包含 47 种身份和凭据模式，以及一个要求移除所有预埋泄漏的预埋泄漏自测；二者报告的命中数均为零。

\subsection{微调配置}\label{app:ft}
我们使用参数高效的监督微调对 Qwen3.8-27B\citep{qwen2026}进行微调：在 4 比特量化的基础权重\citep{dettmers2023}上使用 LoRA\citep{hu2022lora}。每条智能体轨迹被视为一个长训练文档，使模型能够在原始开发上下文中学习规划、编码、工具使用和改进决策。除非另有说明，所有训练配置使用相同的超参数；只有训练语料的组成不同。

\paragraph{训练数据混合。}这三个变体并不是三种独立的混合，而是三个嵌套的语料：每个都在前一个的基础上加入一个阶段的监督，因此 $S1 \subset S2 \subset S3$，表~\ref{tab:ablation} 中任意两行之间的差异都可归因于新加入的那个阶段。两种引擎一起训练；表~\ref{tab:mixture} 给出了组成。

\begin{table}[htbp]
\caption{\textbf{表~\ref{tab:ablation} 中各变体的训练语料。}每个阶段都向前一个语料添加若干行；“累计”一列即该变体所用的训练数据。保留简报列表中的任务在每个阶段都被过滤掉。}
\label{tab:mixture}
\centering\small
\begin{tabular}{@{}llrrrr@{}}
\toprule
新增阶段 & 单位 & Godot & Phaser & 新增行数 & 累计 \\
\midrule
S1 生成 & 智能体轨迹 & 876 & 418 & 1,294 & 1,294 \\
S2 规划 & 简报 $\rightarrow$ 文件规划 & 216 & 92 & 308 & 1,602 \\
S3 改进 & 改进轮次 & 509 & 0 & 509 & 2,111 \\
\midrule
\textit{合计} & & 1,601 & 510 & 2,111 & \\
\bottomrule
\end{tabular}
\end{table}

\paragraph{长上下文训练。}智能体开发轨迹可能长达数万 token。因此，我们使用 57,344 token 的最大上下文长度，并使用 Ulysses 式序列并行\citep{jacobs2023}将每个序列分布到四块 GPU 上。我们使用梯度检查点和融合的交叉熵实现来降低激活值和输出头的显存消耗。

\paragraph{训练目标。}对每条轨迹，原始交互上下文被完整保留，而监督损失只施加于助手生成的输出。这使模型学习开发决策的序列，而不把工具观察或环境反馈当作预测目标。完整的 \rsig{} 模型在规划、生成和经独立验证的改进样本的并集上训练。

\subsection{训练数据消融}\label{app:ablation}
表~\ref{tab:ablation} 给出了第~5 节比较的每个变体（图~\ref{fig:global}c 绘制了 \textsc{+Gen} 和完整模型）。\textsc{Base} 和完整模型分别对应表~\ref{tab:main} 中的 Qwen3.8-27B 和 Qwen3.8-27B (SFT) 行。

\begin{table}[htbp]
\caption{\textbf{用于经验内化的微调配置。}这些是所发布适配器的设置；表~\ref{tab:ablation} 中的每个变体都以此训练，仅在表~\ref{tab:mixture} 的语料上有所不同。}
\label{tab:ftcfg}
\centering\small
\begin{tabular}{@{}ll@{}}
\toprule
配置项 & 取值 \\
\midrule
基础模型 & Qwen3.8-27B \\
适配方式 & LoRA\citep{hu2022lora} \\
LoRA 秩 $r$ & 16 \\
LoRA $\alpha$ & 32 \\
LoRA dropout & 0.05 \\
LoRA 目标模块 & 所有线性投影 \\
秩稳定 LoRA（rsLoRA） & 是 \\
基础权重精度 & 4 比特 NF4\citep{dettmers2023} \\
双重量化 & 是 \\
计算精度 & bfloat16 \\
最大上下文长度 & 57,344 token \\
训练轮数（epoch） & 1 \\
优化器 & AdamW\citep{loshchilov2019}，融合实现 \\
学习率 & $5\times10^{-5}$ \\
学习率调度 & 余弦 \\
预热比例 & 0.03 \\
权重衰减 & 0.01 \\
每设备批大小 & 1 \\
梯度累积 & 无 \\
序列并行度 & 4 \\
硬件 & 4$\times$ H100 80 GB \\
\bottomrule
\end{tabular}
\end{table}

\begin{table}[htbp]
\caption{\textbf{按训练变体划分的 Godot 单次生成结果。}每一行都基于 140 个任务。}
\label{tab:ablation}
\centering\small
\begin{tabular}{@{}lccccc@{}}
\toprule
变体 & 机制 & 深度 & 视觉 & 美术 & 总分 $\uparrow$ \\
\midrule
\textsc{Base} & 41.2 & 33.4 & 38.6 & 38.3 & 37.07 \\
\textsc{+Gen} & 49.0 & 41.4 & 42.0 & 38.0 & 41.45 \\
\textsc{+Gen +Plan} & 48.5 & 40.8 & 42.6 & 39.5 & 41.77 \\
\textsc{+Gen +Plan +Improve} & \textbf{56.1} & \textbf{47.5} & \textbf{49.8} & \textbf{44.8} & \textbf{48.22} \\
\bottomrule
\end{tabular}
\end{table}

\section{扩展评估与案例研究}\label{app:extended}
\subsection{与多智能体“生成–验证”系统的比较}\label{app:vibegame}
\paragraph{为何做这项比较。}本文并不声称智能体能够改进它自己生成的游戏——已经有若干系统做到了这一点。本文声称的是，要\emph{可靠地}做到这一点，需要一个控制器：一个能判断提升空间在哪里、保留最佳版本并适时停止的组件。我们能够运行的最接近的公开系统是 VibeGame\footnote{\url{https://github.com/tettethu/VibeGame}，Apache-2.0。}，这是一个由八个角色组成的智能体团队，其流程本身已包含验证和优化角色。如果这种结构的团队能够自行弥合差距，那么控制器就只是一个装饰。

\paragraph{设置。}使用 GameCraft-Bench 的四十个 Phaser 任务，VibeGame 代码不作修改，其全部八个角色都设为 GPT-5.5，每个任务一次试验，全自动、无人工干预。构建使用与本文所有其他 Phaser 分支相同的验证器、细则和评审打分。

\paragraph{结果说明了什么。}一个本身已经会验证和优化的智能体团队得分为 30.5，而在相同任务上单次生成为 46.0：它的循环并没有把轮次转化为质量。在我们共有的十二个任务上，它达到 23.3，而我们 Phaser 分支的冻结 $P_0$ 已经是 50.1，\rsig{} 将这十二个任务提升到 62.4。流程中存在验证和优化角色，并不是带来持续改进的原因。

\begin{table}[htbp]
\caption{\textbf{VibeGame 与单次生成及 \rsig{} 的对比。}左：VibeGame 的四十个任务与我们四十任务 Phaser 开发子集共有的十二个任务，因此每一列都是同样的十二个游戏。右：VibeGame 自己的四十个任务，与其发布版本在这些任务上引用的单次生成参考结果对比。}
\label{tab:vibegame}
\centering\small
\begin{tabular}{@{}lcc@{}}
\toprule
& 共有的 12 个任务 & VibeGame 自己的 40 个任务 \\
\midrule
VibeGame & 23.3 & 30.5 \\
OpenGame + GPT-5.5，单次生成 & 50.1 & 46.0 \\
\quad + Play2Code & 56.4 & — \\
\rowcolor{rowlight}\quad + \rsig & \textbf{62.4} & — \\
\bottomrule
\end{tabular}
\end{table}

\paragraph{公平地解读。}每个任务只有一次试验，而且我们对该系统的配置仅限于设置其角色模型；对于 VibeGame 在其他模型或更大预算下的表现，我们不作任何论断。

\subsection{定性演化案例研究}\label{app:cases}
我们跟踪三个游戏从代码模型生成的构建到我们最终交付的构建的全过程。每个案例以一条时间线开头：30 轮自主运行中每三轮的官方得分、全局质量监控器持有的版本、饱和停止结束运行的轮次，以及在坐标轴断开之后、循环本身没有运行的那些阶段——冠军不再被超越后写下的阶段简报，以及其后的开发阶段。时间线上的每个构建随后都有一张自己的卡片：该轮或该阶段发现了什么、改变了什么，以及据以评判它的画面；截图上的框是目测放置的，绿色表示某样东西已经到位，红色表示仍然有误。同一个游戏的每个构建都由同一脚本驱动，因此其面板可以相互比较。

\subsubsection{在退化中守住最佳版本}\label{app:case1}
\begin{figure}[h]
\centering
\fig{12}
\caption{\textit{Lawn Guardians}，从生成的游戏到我们交付的构建。得分带给出 30 轮自主运行中每三轮的官方得分、全局质量监控器持有的版本，以及饱和停止在第 21 轮（提前九轮）交付该版本的星标位置。轮次带为每一轮给出一个单元格——改进轮或美术轮——其上方是改进后的验证结论。坐标轴断开之后是循环没有运行的阶段：冠军经受住三个检查点后写下的阶段简报，以及其后的三次迭代；它们是开发阶段而不是循环轮次，因此没有轮次编号。下方的每个构建都在图~\ref{fig:lg1}–\ref{fig:lg2} 中展开为一张卡片。}
\label{fig:lg-timeline}
\end{figure}

开发并不在运行结束的地方结束。全局质量监控器的停止准则关注的是运行自身的得分，而不是游戏本身：一旦冠军经受住三个检查点，循环的证据就不再指向任何地方，而一份从外部写下的阶段简报提供了循环已经失去的方向。其后又有三个阶段——操控、动画，以及配有原创配乐的场景布置——每一个改变的都是游戏玩起来的感受，而不是它做了什么。六个构建在图~\ref{fig:lg1}–\ref{fig:lg2} 中各以一张卡片展示。

\begin{figure}[p]
\centering
\fig{13}
\caption{\textit{Lawn Guardians} 从生成的构建到阶段简报。每张卡片给出该轮或该阶段发现了什么、改变了什么，以及据以评判它的画面；截图上的框为目测放置，绿色表示某样东西已经到位，红色表示仍然有误，每个构建都由同一脚本驱动，因此场景可以相互比较。\textbf{基础游戏：}生成器把规则做对了、把呈现做错了——它附带的美术躺在磁盘上无人使用。\textbf{自主开发：}十二轮把这些美术接入并加上命中反馈，由此产生的冠军再未被超越。\textbf{阶段 1：}冠军被保持了三个检查点后，循环已经用尽了它自己的证据，而一份从外部写下的简报——由 Claude Opus 5 与人类各写一份——让阳光落在草坪上、让僵尸待在各自的行内。}
\label{fig:lg1}
\end{figure}

\begin{figure}[p]
\centering
\fig{14}
\caption{阶段简报之后的三个开发阶段。\textbf{迭代 1} 是一项任何棋盘静止画面都无法展示的修复，因此改为用同一脚本下的两组前后对比来记录：在拉伸的画布上，点击被变换了两次，使植物落在离光标一格远的地方，并且豌豆命中后仍留在草坪上。\textbf{迭代 2} 替换了标题画面和行走循环，是唯一一个静止画面确实能承载的步骤。\textbf{迭代 3} 是场景布置和重写的配乐；种子栏和边缘出现在画面中，而音频则是任何图都无法展示的那一半。}
\label{fig:lg2}
\end{figure}

在 \textit{Lawn Guardians}（图~\ref{fig:lg-timeline}）中，进展并非单调：一个美术轮（R6）和一个反馈轮（R18）各自降低了得分，而全局质量监控器两者都没有采纳。它一直持有 R12（86.1，起点为 61.7），直到饱和停止在第 21 轮交付它。这里也显示出两个局限：验证回放可能太短而看不到真正的修复（R26），而官方评审可能会奖励一个美术遮挡了游戏区域的构建（R27）。Play2Code 的峰值为 81.5，最终为 65.7。

\subsubsection{同一检查点上的两份简报}\label{app:case2}
\textit{Alley Brawlers} 是这样一个案例：循环依据自身证据换来了布景，却失去了演员。三十轮自主开发在打斗背后画出了一条灯火街道——这是整个运行中最大的单项变化——而在同样这些轮次中，格斗者不再能被正确渲染：在交付的构建中，运动中的身体是一团撕裂的涂抹。二者相互抵消，饱和停止交付的冠军得分与基础游戏完全相同。

循环无法得出的观察是：没有人能赢。两位阶段 1 导演都玩了这个构建并指出了这一点，而按模型简报进行的后续运行第一次让得分发生了变化。一个从同一检查点重启、完全不给简报的对照分支达到 60.0；两份简报分别达到 61.5 和 62.0。随后，在阶段 1 冠军上写下的第二份简报针对的是仍然让比赛无法进行的问题——两个身体站在同一空间里——换来的是一场可读的交锋，而不是更高的分数。

\begin{figure}[h]
\centering
\fig{15}
\caption{\textit{Alley Brawlers}，从生成的游戏到我们交付的构建。得分带、轮次带和画面同图~\ref{fig:lg-timeline}。饱和停止在第 18 轮结束自主运行并交付第 9 轮。坐标轴断开之后是循环没有运行的阶段：两份阶段简报（各由一位亲自玩过当时构建的导演写成）以及三个开发阶段。阶段 2 和各迭代都从阶段 1 冠军继续，因此它们是并列的备选，而不是一个序列。}
\label{fig:ab-timeline}
\end{figure}

\begin{figure}[p]
\centering
\fig{16}
\caption{\textit{Alley Brawlers} 从生成的构建到第一份简报。卡片、框和脚本同图~\ref{fig:lg1}。\textbf{基础游戏：}在下着方块雨的扁平条带上进行的一场能运行的比赛。\textbf{自主开发：}场地出现了，格斗者却坏了，交付的构建得分与基础游戏完全相同。\textbf{阶段 1：}一位打了两场比赛的导演要求公平的交锋，以及始终保持为一个完整人物的格斗者；得分在整个运行中第一次发生变化。}
\label{fig:ab1}
\end{figure}

\begin{figure}[p]
\centering
\fig{17}
\caption{第二份简报和各开发阶段，都从阶段 1 冠军继续。\textbf{阶段 2：}第二位导演指出间距是玩家无法行动的原因，交锋由此从双方都变得清晰可读。\textbf{迭代 1–2} 重建了前端——标题从一堵文字墙变成分步流程，HUD 为两名格斗者署名并展示其形象——并加入命中火花和伤害数字。\textbf{迭代 3} 关乎接触的手感：击中停顿、震屏、K.O. 慢动作、连击计数器和配乐。}
\label{fig:ab2}
\end{figure}

\subsubsection{一个循环无法改进的游戏}\label{app:case3}
\textit{Block Cascade} 是相反的起点。七袋随机器、后三块预览、幽灵方块、暂存交换、踢墙、消行、顶出判负和最高分在生成的构建中全都能正常工作，而三十轮由证据引导的改进找不到任何值得替换的版本。循环并不是卡在某个缺陷上；它已经找不到任何可以称为缺陷的东西了。

从该构建重新开始才带来了回报，而带来回报的并不是简报：两个由简报引导的分支和不给简报的对照分支都达到了 50.0，简报换来的是在第 6 轮而非第 12 轮到达那里。第二份简报随后要求增加紧张感——一个有名称的模式、一个实时计时器、一个奖励阶梯——并且得到了它们，同时也带来了新读数与屏幕上已有读数之间的冲突，这由其后的阶段重新布局解决。

\begin{figure}[h]
\centering
\fig{18}
\caption{\textit{Block Cascade}，从生成的游戏到我们交付的构建。生成的游戏已经实现了规格中的每个机制，因此改进轮几乎无事可改：在九个检查点中，全局质量监控器从未偏好过后来的构建，饱和停止在第 9 轮结束运行，而它交付的正是第 0 轮本身。坐标轴断开之后的一切都源自从该构建出发的重启。}
\label{fig:bc-timeline}
\end{figure}

\begin{figure}[p]
\centering
\fig{19}
\caption{\textit{Block Cascade} 从生成的构建经历两份简报。卡片、框和脚本同图~\ref{fig:lg1}。\textbf{基础游戏：}规格中的每个机制第一次就能工作，九轮之后循环交付的仍是这个构建。\textbf{阶段 1：}游戏井有了手绘背景、多面体方块和危险线——而一个不给简报的对照分支达到了同样的得分。\textbf{阶段 2：}有名称的模式和实时速度读数出现了，却被画在了已有读数之上。}
\label{fig:bc1}
\end{figure}

\begin{figure}[t]
\centering
\fig{20}
\caption{结束整个运行的两个开发阶段。\textbf{迭代 1} 为每个数字安排了独立的位置，并为落地加上反馈。\textbf{迭代 2} 增加了变暗棋盘上方的居中结算卡片、一个写明用哪个键填充的 HOLD 框，以及取代 1.8 秒片段的原创配乐。}
\label{fig:bc2}
\end{figure}

\subsection{按类别的结果}\label{app:family}
\newcommand{\famtab}[4]{%
\begin{table}[htbp]
\caption{#1}\label{#2}
\centering\small
\begin{tabular}{@{}lrccc@{}}
\toprule
类别 & $N$ & 基线 $\uparrow$ & #3 \\
\midrule
#4
\bottomrule
\end{tabular}
\end{table}}

\famtab{\textbf{Godot 上 Codex + GPT-5.5 (high) 生成器的按类别细分}（表~\ref{tab:main} 第一个行组），指标为平均总分（$\uparrow$），评审为 Qwen3.8-27B。$N$ 为该类别中的任务数（打分进行中时为“已打分/总数”）；各类别按基线得分排序；未生成项目的任务记 0 分。缩写：+P2C 即 Play2Code；–：尚未在全部任务上可得。}{tab:fam1}{+P2C & +\rsig}{
idle & 4 & 81.9 & 86.8 & 87.5 \\
sports & 4 & 72.8 & 70.5 & 81.2 \\
horror & 5 & 67.2 & 71.4 & 75.1 \\
openworld & 15 & 59.2 & 57.9 & 62.2 \\
tycoon & 16 & 53.3 & 54.4 & 73.2 \\
visualnovel & 11 & 52.0 & 55.8 & 71.5 \\
roguelike & 14 & 46.1 & 38.5 & 56.8 \\
shooter & 7 & 46.0 & 49.7 & 64.1 \\
strategy & 17 & 46.0 & 46.3 & 61.5 \\
rhythm & 5 & 45.4 & 44.7 & 67.9 \\
cardgame & 5 & 44.2 & 44.3 & 71.3 \\
puzzle & 8 & 43.4 & 45.1 & 53.9 \\
platformer & 19 & 43.3 & 44.4 & 58.4 \\
racing & 4 & 43.3 & 53.4 & 68.4 \\
simulation & 6 & 38.1 & 37.7 & 48.8 \\
\midrule
全部 & 140 & 50.26 & 50.74 & 64.53 \\}

\famtab{\textbf{Godot 上 Kimi-K2.6 生成器的按类别细分}（表~\ref{tab:main} 第二个行组）。}{tab:fam2}{+P2C & +\rsig}{
idle & 4 & 53.5 & 50.1 & 60.0 \\
sports & 4 & 45.9 & 57.3 & 64.4 \\
horror & 5 & 43.1 & 56.6 & 63.4 \\
tycoon & 16 & 36.1 & 40.7 & 46.4 \\
rhythm & 5 & 33.7 & 37.1 & 45.4 \\
shooter & 7 & 33.4 & 35.9 & 39.1 \\
simulation & 6 & 30.7 & 29.8 & 34.9 \\
visualnovel & 11 & 30.1 & 36.8 & 45.3 \\
roguelike & 14 & 29.4 & 35.1 & 44.5 \\
racing & 4 & 29.1 & 39.4 & 41.8 \\
openworld & 15 & 27.0 & 33.6 & 42.5 \\
cardgame & 5 & 26.2 & 23.6 & 48.8 \\
strategy & 17 & 24.1 & 33.1 & 47.7 \\
platformer & 19 & 23.3 & 29.2 & 41.4 \\
puzzle & 8 & 15.3 & 19.4 & 29.2 \\
\midrule
全部 & 140 & 29.63 & 35.19 & 44.77 \\}

\famtab{\textbf{Godot 上 GLM-5.3-Flash 生成器的按类别细分}（表~\ref{tab:main} 第三个行组）。}{tab:fam3}{+P2C & +\rsig}{
sports & 4 & 49.6 & 61.5 & 72.6 \\
idle & 4 & 39.0 & 59.2 & 68.6 \\
tycoon & 16 & 38.8 & 43.1 & 56.1 \\
racing & 4 & 34.9 & 36.7 & 53.0 \\
horror & 5 & 34.1 & 43.6 & 49.9 \\
visualnovel & 11 & 32.7 & 38.8 & 59.0 \\
shooter & 7 & 32.6 & 42.5 & 54.9 \\
strategy & 17 & 32.2 & 39.2 & 42.5 \\
simulation & 6 & 32.0 & 38.6 & 46.5 \\
rhythm & 5 & 26.3 & 37.2 & 52.9 \\
cardgame & 5 & 25.5 & 33.1 & 52.7 \\
platformer & 19 & 25.3 & 35.7 & 48.7 \\
roguelike & 14 & 24.9 & 37.5 & 40.9 \\
openworld & 15 & 19.2 & 28.1 & 43.2 \\
puzzle & 8 & 16.9 & 33.3 & 40.8 \\
\midrule
全部 & 140 & 29.46 & 38.59 & 49.72 \\}

\famtab{\textbf{Godot 上 Qwen3.8-27B 生成器的按类别细分}（表~\ref{tab:main} 第四个行组）。}{tab:fam4}{+P2C & +\rsig}{
horror & 5 & 58.7 & 56.6 & 57.0 \\
idle & 4 & 48.8 & 53.2 & 73.9 \\
cardgame & 5 & 48.3 & 53.1 & 59.3 \\
sports & 4 & 44.8 & 51.2 & 55.6 \\
shooter & 7 & 42.3 & 44.7 & 61.5 \\
visualnovel & 11 & 42.0 & 46.0 & 44.9 \\
tycoon & 16 & 41.0 & 52.0 & 62.6 \\
strategy & 17 & 39.5 & 35.2 & 44.9 \\
openworld & 15 & 34.5 & 39.7 & 36.4 \\
rhythm & 5 & 33.6 & 35.0 & 54.7 \\
simulation & 6 & 32.1 & 36.4 & 33.3 \\
platformer & 19 & 31.6 & 31.7 & 43.5 \\
roguelike & 14 & 31.0 & 36.7 & 42.4 \\
puzzle & 8 & 23.5 & 32.5 & 43.5 \\
racing & 4 & 23.3 & 22.7 & 27.9 \\
\midrule
全部 & 140 & 37.07 & 40.53 & 47.77 \\}

\famtab{\textbf{Godot 上 Qwen3.8-27B (SFT) 生成器的按类别细分}（表~\ref{tab:main} 第五个行组）。}{tab:fam5}{Play2Code $\uparrow$ & \rsig{} $\uparrow$}{
idle & 4 & 80.4 & 78.6 & 84.1 \\
sports & 4 & 73.3 & 68.7 & 76.1 \\
horror & 5 & 60.6 & 62.4 & 71.5 \\
tycoon & 16 & 54.7 & 56.4 & 66.3 \\
strategy & 17 & 52.1 & 51.0 & 62.0 \\
visualnovel & 11 & 50.2 & 48.9 & 58.4 \\
openworld & 15 & 48.1 & 46.7 & 58.2 \\
roguelike & 14 & 46.4 & 47.8 & 60.2 \\
platformer & 19 & 42.3 & 46.9 & 56.0 \\
racing & 4 & 41.3 & 38.5 & 52.1 \\
puzzle & 8 & 40.6 & 44.8 & 60.3 \\
simulation & 6 & 40.4 & 45.1 & 61.2 \\
rhythm & 5 & 37.4 & 43.9 & 63.8 \\
cardgame & 5 & 36.1 & 41.8 & 58.9 \\
shooter & 7 & 35.3 & 39.8 & 55.2 \\
\midrule
全部 & 140 & 48.22 & 49.71 & 61.38 \\}

\famtab{\textbf{Phaser 上 OpenGame + GPT-5.5 生成器的按类别细分}（表~\ref{tab:main} Phaser 部分第一个行组）。}{tab:fam6}{+P2C & +\rsig}{
idle & 4 & 72.7 & 73.8 & 78.2 \\
cardgame & 5 & 68.9 & 68.9 & 61.8 \\
visualnovel & 11 & 68.0 & 68.6 & 73.5 \\
tycoon & 16 & 59.9 & 62.9 & 74.3 \\
sports & 4 & 57.2 & 59.8 & 64.8 \\
roguelike & 14 & 56.5 & 58.2 & 66.7 \\
horror & 5 & 53.5 & 52.7 & 57.6 \\
simulation & 6 & 51.8 & 57.7 & 55.4 \\
rhythm & 5 & 50.8 & 52.5 & 58.2 \\
strategy & 17 & 45.7 & 56.8 & 54.8 \\
shooter & 7 & 45.2 & 58.6 & 60.4 \\
racing & 4 & 39.8 & 60.4 & 52.8 \\
platformer & 19 & 38.1 & 43.7 & 45.0 \\
openworld & 15 & 32.8 & 41.0 & 39.5 \\
puzzle & 8 & 32.5 & 40.6 & 51.6 \\
\midrule
全部 & 140 & 49.44 & 55.10 & 58.21 \\}

\famtab{\textbf{Phaser 上 Qwen3.8-27B 生成器的按类别细分}（表~\ref{tab:main} Phaser 部分第二个行组）。}{tab:fam7}{+P2C & +\rsig}{
sports & 4 & 56.1 & 68.8 & 65.4 \\
horror & 5 & 55.6 & 52.9 & 57.6 \\
simulation & 6 & 55.3 & 53.3 & 67.8 \\
visualnovel & 11 & 53.3 & 54.2 & 52.4 \\
rhythm & 5 & 50.9 & 58.4 & 61.4 \\
tycoon & 16 & 47.5 & 58.9 & 56.7 \\
cardgame & 5 & 47.0 & 72.3 & 61.4 \\
puzzle & 8 & 42.6 & 44.2 & 39.5 \\
idle & 4 & 41.3 & 61.5 & 65.2 \\
racing & 4 & 40.9 & 45.3 & 51.6 \\
strategy & 17 & 37.3 & 50.0 & 50.6 \\
platformer & 19 & 37.1 & 40.2 & 46.2 \\
roguelike & 14 & 29.6 & 40.7 & 48.8 \\
shooter & 7 & 27.7 & 46.0 & 48.1 \\
openworld & 15 & 21.1 & 32.0 & 29.3 \\
\midrule
全部 & 140 & 40.03 & 48.70 & 50.24 \\}

\famtab{\textbf{Phaser 上 Qwen3.8-27B (SFT) 生成器的按类别细分}（表~\ref{tab:main} Phaser 部分第三个行组）。}{tab:fam8}{Play2Code $\uparrow$ & \rsig{} $\uparrow$}{
idle & 4 & 73.0 & 75.2 & 78.5 \\
horror & 5 & 56.0 & 60.7 & 67.5 \\
tycoon & 16 & 53.3 & 62.0 & 66.5 \\
sports & 4 & 51.6 & 49.9 & 56.8 \\
shooter & 7 & 41.8 & 51.4 & 55.4 \\
rhythm & 5 & 46.8 & 61.3 & 66.2 \\
simulation & 6 & 46.0 & 50.7 & 60.1 \\
roguelike & 14 & 45.5 & 56.8 & 60.2 \\
openworld & 15 & 42.2 & 46.0 & 50.7 \\
puzzle & 8 & 43.8 & 58.7 & 64.8 \\
strategy & 17 & 42.9 & 51.3 & 54.9 \\
cardgame & 5 & 41.6 & 45.1 & 58.7 \\
visualnovel & 11 & 40.0 & 46.2 & 54.3 \\
platformer & 19 & 39.9 & 50.4 & 54.1 \\
racing & 4 & 32.8 & 39.1 & 49.6 \\
\midrule
全部 & 140 & 45.14 & 53.15 & 58.53 \\}

\end{document}
